% Generated by the blueprint tool (blueprint-pdf/1) from a blueprint snapshot; edit the
% Markdown sources, not this file.  Compile with `latexmk -pdf` or by running
% `pdflatex` twice (the second run settles the contents and the page references).
% Images, if any, are read from assets/ next to this file.
\documentclass[11pt]{article}
\usepackage[T1]{fontenc}
\usepackage[utf8]{inputenc}
\usepackage{lmodern}
\usepackage{textcomp}
\usepackage[a4paper,margin=2.5cm,headheight=14pt]{geometry}
\usepackage{amsmath,amssymb,amsthm}
\usepackage{amscd}  % KaTeX's CD diagrams; part of amsmath, so always there
\usepackage{mathtools}
\usepackage{bm}
\usepackage{graphicx}
\usepackage{xcolor}
\usepackage{array,booktabs,tabularx}
\usepackage{fancyvrb}
\usepackage[normalem]{ulem}
\usepackage{microtype}
% Optional, and only with their fonts: a small TeX installation has the
% packages without the fonts, which fails at the first formula.
\IfFileExists{mathrsfs.sty}{\IfFileExists{rsfs10.tfm}{\usepackage{mathrsfs}}{}}{}
\IfFileExists{stmaryrd.sty}{\IfFileExists{stmary10.tfm}{\usepackage{stmaryrd}}{}}{}
\IfFileExists{cancel.sty}{\usepackage{cancel}}{}
\usepackage[hyphens]{url}
\usepackage[pdfusetitle,bookmarksnumbered=false]{hyperref}
\usepackage{bookmark}
\definecolor{bplink}{HTML}{1F4E9C}
\definecolor{bpmuted}{HTML}{5F676E}
\hypersetup{colorlinks=true,linkcolor=bplink,urlcolor=bplink,citecolor=bplink}
\setlength{\parindent}{0pt}
\setlength{\parskip}{0.55em plus 0.2em minus 0.1em}
% Lists as tight as the page's.
\makeatletter
\def\@listi{\leftmargin\leftmargini\parsep 0pt\topsep 0.3em plus 0.1em\itemsep 0.1em plus 0.1em}
\let\@listI\@listi
\def\@listii{\leftmargin\leftmarginii\labelwidth\leftmarginii\advance\labelwidth-\labelsep\parsep 0pt\topsep 0.1em\itemsep 0.1em}
% An ordered list that does not start at 1.
\newcommand{\bpstartat}[1]{\setcounter{\@enumctr}{\numexpr#1-1\relax}}
\makeatother
\setcounter{tocdepth}{3}
\emergencystretch=3em
\allowdisplaybreaks

% A numbered statement, set as a paper sets one: "Theorem 2.1 (Title)."
% #1 is the lead ("Theorem~2.1"), #2 what follows it in roman.
\newenvironment{bpstatement}[2]{\par\addvspace{0.6\baselineskip}\noindent\textbf{#1}#2\textbf{.}\enspace\ignorespaces}
  {\par\addvspace{0.4\baselineskip}}
\newenvironment{bpstatementit}[2]{\begin{bpstatement}{#1}{#2}\itshape\ignorespaces}{\end{bpstatement}}
% A step: a relation with prose of its own, read under the object it starts from.
\newenvironment{bpstep}[1]{\par\addvspace{0.3\baselineskip}\begin{list}{}{\leftmargin=1.6em\rightmargin=0pt\topsep=0pt}\item\relax
  {\small\color{bpmuted}#1}\par\nopagebreak}{\end{list}\par}
% Notes about where an object is written out.
\newenvironment{bpnote}{\par\small\color{bpmuted}\itshape}{\par}
\newcommand{\bpomitted}[1]{\par{\small\color{bpmuted}\itshape #1\par}}
\newcommand{\bpbodyheading}[1]{\par\addvspace{0.4\baselineskip}\noindent\textbf{#1}\par\nopagebreak}
\newcommand{\bprule}{\par\medskip\noindent\hfil\rule{0.3\linewidth}{0.4pt}\hfil\par\medskip}
\newcommand{\bpboundary}[1]{\par\nopagebreak{\small\color{bpmuted}#1\par}\nopagebreak}
% A reference to something this export leaves out, and one that resolves nowhere.
\newcommand{\bpout}[1]{#1\textsuperscript{\hyperlink{bp-outside}{\ensuremath{\ast}}}}
\newcommand{\bpbroken}[1]{\textcolor{red!70!black}{[#1?]}}
\newcommand{\bpslug}[1]{\texttt{#1}}
% A bibliography entry: label, anchor, entry; the label hangs in the margin.
\newcommand{\bpbibitem}[3]{\par\setbox0\hbox{[#1]\enspace}\ifdim\wd0<5em\setbox0\hbox to 5em{[#1]\hfil}\fi
  \hangindent=\wd0 \hangafter=1 \noindent\box0\phantomsection\label{#2}#3\par\smallskip}
% A character the export knows no way to set.
\newcommand{\bpglyph}[1]{\fbox{\scriptsize #1}}
% Images, at their own size unless they are wider than the line.
\newsavebox{\bpimagebox}
\newcommand{\bpimage}[1]{\sbox{\bpimagebox}{\includegraphics{#1}}%
  \ifdim\wd\bpimagebox>\linewidth\resizebox{\linewidth}{!}{\usebox{\bpimagebox}}\else\usebox{\bpimagebox}\fi}
\newcommand{\bpmissingimage}[1]{\fbox{\parbox{0.8\linewidth}{\centering\small\color{bpmuted}Image not included: #1}}}
\DefineVerbatimEnvironment{bpverbatim}{Verbatim}{fontsize=\small,xleftmargin=1em}
\renewcommand{\qedsymbol}{\ensuremath{\blacksquare}}
\hypersetup{pdftitle={Weil Conjectures — 5 Construction of f!},pdfsubject={Chapter 5 Construction of f ! ; all levels},pdfcreator={blueprint}}
% KaTeX commands the website knows
\providecommand{\R}{\mathbb{R}}\providecommand{\N}{\mathbb{N}}\providecommand{\Z}{\mathbb{Z}}
\providecommand{\Q}{\mathbb{Q}}\providecommand{\C}{\mathbb{C}}
\providecommand{\Reals}{\mathbb{R}}\providecommand{\reals}{\mathbb{R}}\providecommand{\Complex}{\mathbb{C}}
\providecommand{\cnums}{\mathbb{C}}\providecommand{\natnums}{\mathbb{N}}\providecommand{\Bbb}[1]{\mathbb{#1}}
\providecommand{\bold}[1]{\mathbf{#1}}\providecommand{\frak}[1]{\mathfrak{#1}}\providecommand{\mathbbm}[1]{\mathbb{#1}}
\providecommand{\lang}{\langle}\providecommand{\rang}{\rangle}\providecommand{\bra}[1]{\langle #1\rvert}
\providecommand{\ket}[1]{\lvert #1\rangle}\providecommand{\braket}[1]{\langle #1\rangle}
\providecommand{\set}[1]{\{#1\}}\providecommand{\Set}[1]{\left\{#1\right\}}
\providecommand{\infin}{\infty}\providecommand{\isin}{\in}\providecommand{\plusmn}{\pm}\providecommand{\sub}{\subset}
\providecommand{\sube}{\subseteq}\providecommand{\supe}{\supseteq}\providecommand{\exist}{\exists}
\providecommand{\bull}{\bullet}\providecommand{\larr}{\leftarrow}\providecommand{\rarr}{\rightarrow}
\providecommand{\lrarr}{\leftrightarrow}\providecommand{\Larr}{\Leftarrow}\providecommand{\Rarr}{\Rightarrow}
\providecommand{\Lrarr}{\Leftrightarrow}\providecommand{\uarr}{\uparrow}\providecommand{\darr}{\downarrow}
\providecommand{\harr}{\leftrightarrow}\providecommand{\lArr}{\Leftarrow}\providecommand{\rArr}{\Rightarrow}
\providecommand{\hArr}{\Leftrightarrow}\providecommand{\alef}{\aleph}\providecommand{\alefsym}{\aleph}
\providecommand{\real}{\Re}\providecommand{\image}{\Im}\providecommand{\weierp}{\wp}
\providecommand{\Alpha}{\mathrm{A}}\providecommand{\Beta}{\mathrm{B}}\providecommand{\Epsilon}{\mathrm{E}}
\providecommand{\Zeta}{\mathrm{Z}}\providecommand{\Eta}{\mathrm{H}}\providecommand{\Iota}{\mathrm{I}}
\providecommand{\Kappa}{\mathrm{K}}\providecommand{\Mu}{\mathrm{M}}\providecommand{\Nu}{\mathrm{N}}
\providecommand{\Omicron}{\mathrm{O}}\providecommand{\Rho}{\mathrm{P}}\providecommand{\Tau}{\mathrm{T}}
\providecommand{\Chi}{\mathrm{X}}\providecommand{\omicron}{o}\providecommand{\gt}{>}\providecommand{\lt}{<}
\providecommand{\And}{\&}\providecommand{\hdashline}{\hline}\providecommand{\operatornamewithlimits}[1]{\operatorname*{#1}}
\providecommand{\llbracket}{[\![}\providecommand{\rrbracket}{]\!]}\providecommand{\cancel}[1]{#1}
\providecommand{\mathscr}{\mathcal}
% the characters beyond ASCII this document uses
\DeclareUnicodeCharacter{2160}{\textup{I}}
\DeclareUnicodeCharacter{2161}{\textup{II}}
% the project's macros ([katex.macros] in blueprint.toml)
\providecommand{\A}{}\renewcommand{\A}{\mathbf A}
\providecommand{\AEn}{}\renewcommand{\AEn}{A_{E,n}}
\providecommand{\Aaf}{}\renewcommand{\Aaf}{\mathbf A}
\providecommand{\Ab}{}\renewcommand{\Ab}{\mathsf{Ab}}
\providecommand{\AffSch}{}\renewcommand{\AffSch}{\operatorname{AffSch}}
\providecommand{\Aone}{}\renewcommand{\Aone}{\mathbf A^1}
\providecommand{\Bl}{}\renewcommand{\Bl}{\operatorname{Bl}}
\providecommand{\Br}{}\renewcommand{\Br}{\mathrm{Br}}
\providecommand{\C}{}\renewcommand{\C}{\mathbb C}
\providecommand{\Cech}{}\renewcommand{\Cech}{\check{H}}
\providecommand{\Coext}{}\renewcommand{\Coext}{\operatorname{Coext}}
\providecommand{\Cone}{}\renewcommand{\Cone}{\operatorname{Cone}}
\providecommand{\Cons}{}\renewcommand{\Cons}{\operatorname{Cons}}
\providecommand{\Cref}{}\renewcommand{\Cref}[1]{\text{#1}}
\providecommand{\D}{}\renewcommand{\D}{D}
\providecommand{\DLAFT}{}\renewcommand{\DLAFT}{\operatorname{FT}}
\providecommand{\DLAFr}{}\renewcommand{\DLAFr}{\operatorname{Fr}}
\providecommand{\DLAQ}{}\renewcommand{\DLAQ}{{\overline{\mathbf Q}_\ell}}
\providecommand{\DLRQ}{}\renewcommand{\DLRQ}{{\overline{\mathbf Q}_\ell}}
\providecommand{\Dbc}{}\renewcommand{\Dbc}{D^b_c}
\providecommand{\Dc}{}\renewcommand{\Dc}{D_{\mathrm{cons}}}
\providecommand{\Dctf}{}\renewcommand{\Dctf}{D^b_{ctf}}
\providecommand{\Dic}{}\renewcommand{\Dic}{D_{\IndCons}}
\providecommand{\Dplus}{}\renewcommand{\Dplus}{D^{+}}
\providecommand{\End}{}\renewcommand{\End}{\operatorname{End}}
\providecommand{\Et}{}\renewcommand{\Et}{\mathrm{\acute{e}t}}
\providecommand{\Ex}{}\renewcommand{\Ex}{\mathrm{Ex}}
\providecommand{\Ext}{}\renewcommand{\Ext}{\operatorname{Ext}}
\providecommand{\F}{}\renewcommand{\F}{\mathbb F}
\providecommand{\FT}{}\renewcommand{\FT}{\operatorname{FT}}
\providecommand{\Fbar}{}\renewcommand{\Fbar}{{\overline{\mathbf F}_q}}
\providecommand{\Fell}{}\renewcommand{\Fell}{\mathbf F_\ell}
\providecommand{\Fib}{}\renewcommand{\Fib}{\operatorname{Fib}}
\providecommand{\Fix}{}\renewcommand{\Fix}{\operatorname{Fix}}
\providecommand{\Fp}{}\renewcommand{\Fp}{{\mathbf F_{p}}}
\providecommand{\Fq}{}\renewcommand{\Fq}{{\mathbf F_q}}
\providecommand{\Fqn}{}\renewcommand{\Fqn}{{\mathbf F_{q^{n}}}}
\providecommand{\Frob}{}\renewcommand{\Frob}{\operatorname{Frob}}
\providecommand{\GL}{}\renewcommand{\GL}{\operatorname{GL}}
\providecommand{\Gal}{}\renewcommand{\Gal}{\operatorname{Gal}}
\providecommand{\Hens}{}\renewcommand{\Hens}{\operatorname{Hens}}
\providecommand{\Hilb}{}\renewcommand{\Hilb}{\operatorname{Hilb}}
\providecommand{\Hom}{}\renewcommand{\Hom}{\operatorname{Hom}}
\providecommand{\IndCons}{}\renewcommand{\IndCons}{\operatorname{IndCons}}
\providecommand{\Kar}{}\renewcommand{\Kar}{\operatorname{Kar}}
\providecommand{\LT}{}\renewcommand{\LT}{\operatorname{LT}}
\providecommand{\Lie}{}\renewcommand{\Lie}{\operatorname{Lie}}
\providecommand{\Loc}{}\renewcommand{\Loc}{\operatorname{Loc}}
\providecommand{\Mod}{}\renewcommand{\Mod}{\mathsf{Mod}}
\providecommand{\N}{}\renewcommand{\N}{\mathbb{N}}
\providecommand{\Needspace}{}\renewcommand{\Needspace}{\par}
\providecommand{\O}{}\renewcommand{\O}{{\mathcal{O}}}
\providecommand{\Ob}{}\renewcommand{\Ob}{\mathrm{Ob}}
\providecommand{\Oe}{}\renewcommand{\Oe}{\mathcal O_E}
\providecommand{\PP}{}\renewcommand{\PP}{\mathbf{P}}
\providecommand{\Perf}{}\renewcommand{\Perf}{\operatorname{Perf}}
\providecommand{\Pic}{}\renewcommand{\Pic}{\mathrm{Pic}}
\providecommand{\Pm}{}\renewcommand{\Pm}{\mathbf P^{m}}
\providecommand{\Pone}{}\renewcommand{\Pone}{\mathbf P^{1}}
\providecommand{\Pqn}{}\renewcommand{\Pqn}{\mathbf P^{n}}
\providecommand{\ProFin}{}\renewcommand{\ProFin}{\operatorname{ProFin}}
\providecommand{\Proet}{}\renewcommand{\Proet}{\mathrm{pro\acute{e}t}}
\providecommand{\Q}{}\renewcommand{\Q}{\mathbb Q}
\providecommand{\Qb}{}\renewcommand{\Qb}{\overline{\mathbf Q}_{\ell}}
\providecommand{\Qbar}{}\renewcommand{\Qbar}{\overline{\mathbf Q}}
\providecommand{\QbarX}{}\renewcommand{\QbarX}{\overline{\mathbf Q}_{\ell,X}}
\providecommand{\Qbarl}{}\renewcommand{\Qbarl}{{\overline{\mathbf Q}_{\ell}}}
\providecommand{\Qell}{}\renewcommand{\Qell}{{\mathbf{Q}_{\ell}}}
\providecommand{\Ql}{}\renewcommand{\Ql}{\mathbf Q_{\ell}}
\providecommand{\Qlb}{}\renewcommand{\Qlb}{{\overline{\mathbf Q}_\ell}}
\providecommand{\Qlbar}{}\renewcommand{\Qlbar}{\overline{\mathbf Q}_\ell}
\providecommand{\R}{}\renewcommand{\R}{\mathbb{R}}
\providecommand{\RG}{}\renewcommand{\RG}{\operatorname{\mathrm{R}\Gamma}}
\providecommand{\RGamma}{}\renewcommand{\RGamma}{R\Gamma}
\providecommand{\RGammac}{}\renewcommand{\RGammac}{R\Gamma_c}
\providecommand{\RHom}{}\renewcommand{\RHom}{R\mathcal{H}om}
\providecommand{\RHomX}{}\renewcommand{\RHomX}{R\mathcal H om_X}
\providecommand{\Rat}{}\renewcommand{\Rat}{\operatorname{IsRat}}
\providecommand{\Res}{}\renewcommand{\Res}{\operatorname{Res}}
\providecommand{\Rlim}{}\renewcommand{\Rlim}{R\!\varprojlim}
\providecommand{\SIm}{}\renewcommand{\SIm}{\operatorname{SIm}}
\providecommand{\Sch}{}\renewcommand{\Sch}{\mathsf{Sch}}
\providecommand{\Sh}{}\renewcommand{\Sh}{\operatorname{Sh}}
\providecommand{\Span}{}\renewcommand{\Span}{\operatorname{span}}
\providecommand{\Spec}{}\renewcommand{\Spec}{\operatorname{Spec}}
\providecommand{\Tr}{}\renewcommand{\Tr}{\operatorname{Tr}}
\providecommand{\Trsm}{}\renewcommand{\Trsm}{\operatorname{Tr}^{\mathrm{sm}}}
\providecommand{\Wt}{}\renewcommand{\Wt}{\operatorname{Wt}}
\providecommand{\Z}{}\renewcommand{\Z}{\mathbb{Z}}
\providecommand{\Zar}{}\renewcommand{\Zar}{\mathrm{zar}}
\providecommand{\Zell}{}\renewcommand{\Zell}{{\mathbf{Z}_{\ell}}}
\providecommand{\Zl}{}\renewcommand{\Zl}{{\mathbf Z_\ell}}
\providecommand{\affet}{}\renewcommand{\affet}[1]{{#1}^{\mathrm{aff}}_{\mathrm{ét}}}
\providecommand{\affproet}{}\renewcommand{\affproet}[1]{{#1}^{\mathrm{aff}}_{\mathrm{proét}}}
\providecommand{\boxt}{}\renewcommand{\boxt}{\boxtimes_A^{\mathrm L}}
\providecommand{\cF}{}\renewcommand{\cF}{\mathcal{F}}
\providecommand{\cHom}{}\renewcommand{\cHom}{R\mathcal{H}\!om}
\providecommand{\cI}{}\renewcommand{\cI}{\mathcal{I}}
\providecommand{\cL}{}\renewcommand{\cL}{\mathcal L}
\providecommand{\cM}{}\renewcommand{\cM}{\mathcal M}
\providecommand{\cN}{}\renewcommand{\cN}{\mathcal N}
\providecommand{\cO}{}\renewcommand{\cO}{\mathcal O}
\providecommand{\cT}{}\renewcommand{\cT}{\mathcal T}
\providecommand{\checked}{}\renewcommand{\checked}{\color{red}\aftergroup\restorecolor}
\providecommand{\coeff}{}\renewcommand{\coeff}{\operatorname{coeff}}
\providecommand{\coker}{}\renewcommand{\coker}{\operatorname{coker}}
\providecommand{\colim}{}\renewcommand{\colim}{\operatorname{colim}}
\providecommand{\constructedby}{}\renewcommand{\constructedby}[1]{\par\smallskip\noindent\textcolor{purple!60!black}{\textsf{Constructed by:} }
  \footnotesize\sloppy\nolinkurl{#1}\par}
\providecommand{\deepresult}{}\renewcommand{\deepresult}[1]{\par\smallskip\noindent\textcolor{orange!45!black}{
  \textsf{Blueprint role: named deep theorem node; see the cluster proof architecture.}}
  \par\noindent\textcolor{orange!45!black}{\textsf{Mathematical source:} }#1\par}
\providecommand{\dt}{}\renewcommand{\dt}{\mathbin{\otimes^{\mathbf L}}}
\providecommand{\dw}{}\renewcommand{\dw}{\operatorname{dw}}
\providecommand{\et}{}\renewcommand{\et}[1]{{#1}_{\mathrm{ét}}}
\providecommand{\ev}{}\renewcommand{\ev}{\operatorname{ev}}
\providecommand{\fib}{}\renewcommand{\fib}{\operatorname{fib}}
\providecommand{\formalizationnote}{}\renewcommand{\formalizationnote}[1]{\begin{quote}\small\raggedright\textbf{Formalization note.} #1\end{quote}}
\providecommand{\fpqc}{}\renewcommand{\fpqc}{X_{\mathrm{fpqc}}}
\providecommand{\gammacont}{}\renewcommand{\gammacont}{\Gamma_{\mathrm{cont}}}
\providecommand{\gm}{}\renewcommand{\gm}{\mathbf G_m}
\providecommand{\hocolim}{}\renewcommand{\hocolim}{\operatorname{hocolim}}
\providecommand{\holim}{}\renewcommand{\holim}{\operatorname{holim}}
\providecommand{\hrB}{}\renewcommand{\hrB}{\@ifnextchar]{\hrEnd}{\hrStep}}
\providecommand{\hrEnd}{}\renewcommand{\hrEnd}{\@firstoftwo{}}
\providecommand{\hrStep}{}\renewcommand{\hrStep}{\@firstoftwo{\hrB}}
\providecommand{\id}{}\renewcommand{\id}{\operatorname{id}}
\providecommand{\kunFl}{}\renewcommand{\kunFl}{\mathbf F_\ell}
\providecommand{\kunZn}{}\renewcommand{\kunZn}{\mathbf Z/n\mathbf Z}
\providecommand{\kuncd}{}\renewcommand{\kuncd}{\operatorname{cd}}
\providecommand{\kundim}{}\renewcommand{\kundim}{\dim}
\providecommand{\kunotimes}{}\renewcommand{\kunotimes}[1]{\otimes^{\mathbf L}_{#1}}
\providecommand{\kunsupp}{}\renewcommand{\kunsupp}{\operatorname{Supp}}
\providecommand{\lotimes}{}\renewcommand{\lotimes}{\otimes^{\mathrm L}}
\providecommand{\m}{}\renewcommand{\m}{{\mathfrak{m}}}
\providecommand{\n}{}\renewcommand{\n}{{\mathfrak{n}}}
\providecommand{\nagataGap}{}\renewcommand{\nagataGap}[1]{\par\noindent\textbf{Formalization gap.} #1\par}
\providecommand{\openobligation}{}\renewcommand{\openobligation}[1]{\par\smallskip\noindent\textcolor{red!65!black}{\textsf{Open obligation:} }
  \footnotesize\sloppy\nolinkurl{#1}\par}
\providecommand{\p}{}\renewcommand{\p}{{\mathfrak{p}}}
\providecommand{\pr}{}\renewcommand{\pr}{\operatorname{pr}}
\providecommand{\preshv}{}\renewcommand{\preshv}[1]{\mathrm{PreShv}(#1)}
\providecommand{\pro}{}\renewcommand{\pro}[1]{\mathrm{Pro}(#1)}
\providecommand{\proet}{}\renewcommand{\proet}{X_{\mathrm{proét}}}
\providecommand{\q}{}\renewcommand{\q}{{\mathfrak{q}}}
\providecommand{\rad}{}\renewcommand{\rad}{\operatorname{rad}}
\providecommand{\rank}{}\renewcommand{\rank}{\operatorname{rank}}
\providecommand{\red}{}\renewcommand{\red}{\color{red}}
\providecommand{\restorecolor}{}\renewcommand{\restorecolor}{\color{black}}
\providecommand{\retiredobligation}{}\renewcommand{\retiredobligation}[1]{\par\smallskip\noindent\textcolor{red!75!black}{\textsf{Retired invalid route:} }
  \footnotesize\sloppy\nolinkurl{#1}\par}
\providecommand{\rk}{}\renewcommand{\rk}{\operatorname{rk}}
\providecommand{\sep}{}\renewcommand{\sep}{\textnormal{sep}}
\providecommand{\shv}{}\renewcommand{\shv}[1]{\mathrm{Shv}(#1)}
\providecommand{\spec}{}\renewcommand{\spec}[1]{\mathrm{Spec}(#1)}
\providecommand{\stacks}{}\renewcommand{\stacks}[1]{\href{https://stacks.math.columbia.edu/tag/#1}{\textnormal{Stacks, Tag~#1}}}
\providecommand{\stacksproject}{}\renewcommand{\stacksproject}[1]{{\cite[\href{https://stacks.math.columbia.edu/tag/#1}{Tag #1}]{stacks-project}}}
\providecommand{\statusnote}{}\renewcommand{\statusnote}[1]{\begin{quote}\small\raggedright\textbf{Status.} #1\end{quote}}
\providecommand{\tildeideal}{}\renewcommand{\tildeideal}[2]{\widetilde{(#1, #2)}}
\providecommand{\tildering}{}\renewcommand{\tildering}[3]{{#1}_{\widetilde{(#2, #3)}}}
\providecommand{\tr}{}\renewcommand{\tr}{\operatorname{tr}}
\providecommand{\wt}{}\renewcommand{\wt}{\operatorname{wt}}

\begin{document}
\pagestyle{plain}
\begin{center}
{\LARGE\bfseries Weil Conjectures\par}
\vspace{0.7em}
{\large Chapter 5: Construction of $f^!$\par}
\vspace{0.5em}
{\small\color{bpmuted}Snapshot 2026-09-26, revision 8e4cb88e2a\par}
\end{center}
\vspace{0.5em}
{\small\tableofcontents}
\bigskip\hrule\bigskip

% object: chap-Construction-of-f (blueprint/chap-Construction-of-f/_section.md)
\section*{\texorpdfstring{5\quad Construction of $f^!$}{5 Construction of f!}}
\phantomsection\label{bp1}%
\addcontentsline{toc}{section}{\texorpdfstring{5\quad Construction of $f^!$}{5 Construction of f!}}
% text: prose-05-01
\phantomsection\label{bp2}%
Verdier duality, Poincaré duality and the trace formula all rest on a right adjoint \(f^!\) of \(Rf_!\). Indeed, Verdier duality (Chapter \bpout{10}) is duality with respect to the dualizing complex \(a_X^!\Lambda\), where \(a_X:X\to\operatorname{Spec}k\) is the structure map, and Poincaré duality follows from it once smooth purity (Theorem \bpout{9.3}) computes \(f^!\) for smooth morphisms. Similarly, the local terms of the trace formula (Chapter \bpout{12}) are built from \(f^!\) for the two projections of a correspondence.

The functor \(Rf_!\) itself is easy to write down. A separated morphism of finite type \(f:X\to Y\) has a compactification \(f=p\circ j\), with \(j:X\hookrightarrow\overline X\) an open immersion and \(p:\overline X\to Y\) proper (Theorem \bpout{A.2}), and we put \(Rf_!=Rp_*\circ j_!\). This definition stays within constructible complexes because, for every admissible \(\Lambda\), the functor \(Rp_*\) maps \(D^b_{\mathrm{ctf}}(\overline X,\Lambda)\) into \(D^b_{\mathrm{ctf}}(Y,\Lambda)\). The resulting functor satisfies proper base change and the projection formula (Proposition \bpout{4.40}).

The right adjoint of \(Rf_!\) is another matter. The factor \(j_!\) has the right adjoint \(j^*\), but for \(Rp_*\) with \(p\) proper no formula for a right adjoint is available. Even where a formula exists, it is not computed pointwise: for a closed immersion \(i\) the functor \(i^!\) takes sections with support in the closed subscheme, so the value of \(i^!K\) at a point depends on \(K\) near the point, not only on its stalk there.

Grothendieck and Hartshorne built the coherent \(f^!\) from local formulas and then had to glue them. Neeman noticed, instead, that \(D_{\mathrm{qc}}(X)\) is generated by compact perfect complexes and that \(Rf_*\) preserves direct sums, so that Brown representability produces \(f^!\) with no formula. We follow his approach, with \(D^b_{\mathrm{ctf}}\) in place of the perfect complexes.

The theorem we use is the following. Let \(\mathcal T\) be a triangulated category with arbitrary direct sums, and suppose that \(\mathcal T\) is compactly generated: there is a set \(\mathcal G\) of compact objects \(G\) (that is, \(\Hom(G,-)\) commutes with direct sums) such that \(K=0\) whenever \(\Hom(G[n],K)=0\) for all \(G\in\mathcal G\) and \(n\in\mathbf Z\). Then every exact functor out of \(\mathcal T\) that preserves direct sums has an exact right adjoint [Stacks, \href{https://stacks.math.columbia.edu/tag/0A8F}{0A8F}, \href{https://stacks.math.columbia.edu/tag/0A8G}{0A8G}]. The category \(D^b_{\mathrm{ctf}}(X,\Lambda)\) is not closed under infinite direct sums. We therefore apply the theorem to the localizing category \(\Loc(D^b_{\mathrm{ctf}}(X,\Lambda))\) that it generates (Section \hyperref[bp3]{5.1}). The result is Theorem \hyperref[bp100]{5.31}: for every admissible \(\Lambda\) and every separated morphism of finite type \(f:X\to Y\) there is an exact functor

\[
f^!:\Loc(D^b_{\mathrm{ctf}}(Y,\Lambda))\to\Loc(D^b_{\mathrm{ctf}}(X,\Lambda)),
\qquad
\Hom(Rf_!M,N)\simeq\Hom(M,f^!N).
\]

For finite coefficients nothing is lost by passing to the localizing category, since \(\Loc(D^b_{\mathrm{ctf}}(X,A))\) is equivalent to the classical derived category \(D(X_{\acute{e}t},A)\) (Proposition \hyperref[bp48]{5.15}).

Everything therefore rests on the two hypotheses of the theorem: that the objects of \(D^b_{\mathrm{ctf}}(X,\Lambda)\) are compact in \(\Loc(D^b_{\mathrm{ctf}}(X,\Lambda))\), and that \(Rp_*\) preserves direct sums for proper \(p\). Neither is formal, because the summands of a direct sum may live in degrees going off to \(-\infty\). If the Ext groups out of a sheaf do not vanish in high degrees, compactness fails even for the constant sheaf. Here is an example, outside our conventions because \(\mathbf R\) is neither finite nor algebraically closed. Take \(\ell=2\) and \(X=\operatorname{Spec}\mathbf R\). Then sheaves on \(X_{\acute{e}t}\) are modules over \(\operatorname{Gal}(\mathbf C/\mathbf R)=\mathbf Z/2\), and \(\Ext^n(\mathbf F_2,\mathbf F_2)=H^n(\mathbf Z/2,\mathbf F_2)=\mathbf F_2\) for all \(n\geq0\). Since products of modules are exact and the complexes \(\mathbf F_2[n]\) live in distinct degrees, the sum \(\bigoplus_{n\geq0}\mathbf F_2[n]\) in \(D(X_{\acute{e}t},\mathbf F_2)\) is also the product. Hence

\[
\Hom\Bigl(\mathbf F_2,\bigoplus_{n\ge0}\mathbf F_2[n]\Bigr)
=\prod_{n\ge0}\mathbf F_2
\ne\bigoplus_{n\ge0}\mathbf F_2
=\bigoplus_{n\ge0}\Hom(\mathbf F_2,\mathbf F_2[n]).
\]

By our conventions (Section \bpout{1.9}) the base field is finite or algebraically closed, so its \(\ell\)-cohomological dimension is finite (Lemma \bpout{2.46}). For finite coefficients this rules out such examples. In that case both hypotheses reduce to the fact that classical étale cohomology commutes with direct sums of unbounded complexes (Section \hyperref[bp43]{5.3}). For \(\Lambda=E\) or \(\overline{\mathbf Q}_\ell\) a bound on the base is not enough, and we prove instead a bound on \(\Ext\) that depends only on the source (Lemma \hyperref[bp20]{5.6}).

% object: sec-Existence-and-finiteness-of-f (blueprint/chap-Construction-of-f/sec-Existence-and-finiteness-of-f/_section.md), content hidden
\subsection*{\texorpdfstring{5.1\quad Existence and finiteness of $f^!$}{5.1 Existence and finiteness of f!}}
\phantomsection\label{bp3}%
\addcontentsline{toc}{subsection}{\texorpdfstring{5.1\quad Existence and finiteness of $f^!$}{5.1 Existence and finiteness of f!}}
% text: prose-05-02
\phantomsection\label{bp4}%
We construct \(f^!\) by Brown representability, as outlined at the beginning of this chapter, without writing it down. Brown representability needs a category with arbitrary direct sums, and \(D^b_{\mathrm{ctf}}(X,\Lambda)\) is not closed under them. We therefore enlarge it to the smallest category that has them and show that \(D^b_{\mathrm{ctf}}(X,\Lambda)\) compactly generates it.

% object: definition-localizing-subcategory (blueprint/chap-Construction-of-f/sec-Existence-and-finiteness-of-f/definition-localizing-subcategory/_section.md)
\phantomsection\label{bp5}%
\begin{bpstatement}{Definition~5.1}{ (The localizing subcategory generated by $\Dctf(X,\Lambda)$)}
For admissible \(\Lambda\), let \(\Loc(D^b_{\mathrm{ctf}}(X,\Lambda))\) be the smallest full triangulated subcategory of \(D(X_{\mathrm{pro}\acute{e}t},\Lambda)\) containing \(D^b_{\mathrm{ctf}}(X,\Lambda)\) and closed under arbitrary direct sums.
\end{bpstatement}

% text: prose-05-03
\phantomsection\label{bp6}%
For integral and rational coefficients, this is also \(\Loc(\Dbc(X,\Lambda))\), by Proposition \bpout{3.12}.

For \(\mathcal O_E\) the objects of \(D^b_{\mathrm{ctf}}(X,\mathcal O_E)\) are derived complete, but the localizing subcategory they generate also contains complexes that are not. For this reason the proof of Theorem theorem-exceptional-inverse-image must show, for integral constructible \(K\), that \(f^!K\) is derived complete.

% object: theorem-exceptional-inverse-image (blueprint/chap-Construction-of-f/sec-Existence-and-finiteness-of-f/theorem-exceptional-inverse-image/_section.md)
\phantomsection\label{bp7}%
\begin{bpstatementit}{Theorem~5.2}{ (The exceptional inverse image)}
Let \(\Lambda\in\{A,\mathbf Z_\ell,\mathcal O_E,\mathbf Q_\ell,E,\overline{\mathbf Q}_\ell\}\) be an admissible coefficient ring and \(f:X\to Y\) a morphism of separated schemes of finite type over \(k\).

\begin{enumerate}
\item The objects of \(D^b_{\mathrm{ctf}}(X,\Lambda)\) are compact in \(\Loc(D^b_{\mathrm{ctf}}(X,\Lambda))\) and a small skeleton compactly generates it.
\item Compactifications define an exact functor \(Rf_!:\Loc(D^b_{\mathrm{ctf}}(X,\Lambda))\to\Loc(D^b_{\mathrm{ctf}}(Y,\Lambda))\), preserving direct sums and sending \(D^b_{\mathrm{ctf}}(X,\Lambda)\) into \(D^b_{\mathrm{ctf}}(Y,\Lambda)\).
\item The functor \(Rf_!\) has an exact right adjoint \(f^!:\Loc(D^b_{\mathrm{ctf}}(Y,\Lambda))\to\Loc(D^b_{\mathrm{ctf}}(X,\Lambda))\), obtained from Brown representability.
\item This right adjoint sends \(D^b_{\mathrm{ctf}}(Y,\Lambda)\) into \(D^b_{\mathrm{ctf}}(X,\Lambda)\), so its adjunction restricts to \(D^b_{\mathrm{ctf}}\).
\end{enumerate}
\end{bpstatementit}

% text: prose-05-04
\phantomsection\label{bp8}%
By Proposition \bpout{3.12}, the theorem holds with \(\Dbc\) in place of \(D^b_{\mathrm{ctf}}\) for \(\Lambda=\mathbf Z_\ell,\mathcal O_E,\mathbf Q_\ell,E,\overline{\mathbf Q}_\ell\).

For proper \(p\) we may use the trivial compactification, so that \(Rp_!=Rp_*\). By part (4), the functor \(Rp_*:D^b_{\mathrm{ctf}}(X,\Lambda)\to D^b_{\mathrm{ctf}}(Y,\Lambda)\) then has the right adjoint \(p^!\).

Parts (1)–(3) are proved in this chapter and part (4) in Chapter \bpout{6}. For smooth \(f\), smooth purity (Chapter \bpout{9}) gives a global isomorphism \(f^!\simeq f^*(d)[2d]\), compatible with trace and composition.


% object: sec-Compactness-of-constructible-complexes (blueprint/chap-Construction-of-f/sec-Compactness-of-constructible-complexes/_section.md), content hidden
\subsection*{\texorpdfstring{5.2\quad Compactness of constructible complexes}{5.2 Compactness of constructible complexes}}
\phantomsection\label{bp9}%
\addcontentsline{toc}{subsection}{\texorpdfstring{5.2\quad Compactness of constructible complexes}{5.2 Compactness of constructible complexes}}
% text: prose-05-05
\phantomsection\label{bp10}%
Brown representability produces \(f^!\) once the bounded constructible complexes are known to be compact, that is, once \(\Hom(C,-)\) commutes with direct sums for constructible \(C\). The difficulty is that the summands may have cohomology in degrees going off to \(-\infty\), so that \(\Hom(C,-)\) sees all the groups \(\Ext^n(C,-)\) at once. For instance, over \(\operatorname{Spec}\mathbf R\) these groups never stop, and we saw at the beginning of this chapter that compactness fails there. By contrast, over a finite field they stop, as the following example shows.

% object: example-compactness-finite-field (blueprint/chap-Construction-of-f/sec-Compactness-of-constructible-complexes/example-compactness-finite-field/_section.md)
\phantomsection\label{bp11}%
\begin{bpstatement}{Example~5.3}{ (Compactness over a finite field)}
Let \(X=\operatorname{Spec}\mathbf F_q\) and \(\Lambda=E\). Then \(\Ext^n(E,E)=H^n(X_{\mathrm{pro}\acute{e}t},E)\) is \(E\) for \(n=0,1\) and vanishes for \(n\geq2\). To see the vanishing, consider first \(\mathbf Z_\ell\) and the Milnor sequence (2.11) with \(F_m=(\mathbf Z/\ell^m)_X\). Since the field \(\mathbf F_q\) has \(\ell\)-cohomological dimension one, we have \(H^n(X,\mathbf Z/\ell^m)=0\) for \(n\ge2\) and all \(m\), while the finite groups \(H^1(X,\mathbf Z/\ell^m)\) have vanishing \(\varprojlim^1\). The Milnor sequence therefore gives \(H^n(X_{\mathrm{pro}\acute{e}t},\mathbf Z_{\ell,X})=0\) for \(n\geq2\). From here, inverting \(\ell\) and using Example \bpout{2.57} for \(n=1\) gives the claim for \(\mathbf Q_\ell\); for \(E\neq\mathbf Q_\ell\) we tensor with \(E\) over \(\mathbf Q_\ell\).

As in the example over \(\mathbf R\), products are exact (Proposition \bpout{2.23}) and the complexes \(E[n]\) live in distinct degrees. Hence \(\bigoplus_{n\ge0}E[n]=\prod_{n\ge0}E[n]\), and

\[
\Hom\Bigl(E,\bigoplus_{n\ge0}E[n]\Bigr)
=\prod_{n\ge0}\Ext^n(E,E)
=E\oplus E
=\bigoplus_{n\ge0}\Hom(E,E[n]).
\]

The product equals the sum because only two factors are nonzero. Over \(\mathbf R\) every factor is nonzero, and the product is larger.
\end{bpstatement}

% text: prose-05-06
\phantomsection\label{bp12}%
The general argument rests on the same mechanism: a bound \(d_C\), valid for all constructible targets at once, beyond which \(\Ext^q(C,-)\) vanishes. Until the end of Section \hyperref[bp9]{5.2}, \(\Lambda\) is \(E\) or \(\overline{\mathbf Q}_\ell\), and we show that bounded constructible complexes are compact in \(\Loc(\Dbc(X,\Lambda))\). Since \(D^b_{\mathrm{ctf}}(X,\Lambda)=\Dbc(X,\Lambda)\), these are the proposed compact generators. The advantage of working with \(\Dbc\) is that we can use truncations and the constructible heart, which is abelian by Lemma \bpout{3.18}. As for notation, \(\Hom\) and \(\Ext\) without decoration are taken in \(D(X_{\mathrm{pro}\acute{e}t},\Lambda)\), also when both arguments are constructible sheaves, and a subscript indicates the coefficient ring where it changes, as in the lattice arguments over \(\mathcal O_E\).

By dévissage it suffices to show, for a constructible sheaf \(C\) and a family \((K_\lambda)\) of objects of \(\Loc(\Dbc(X,\Lambda))\), that \(\bigoplus_\lambda\Hom(C,K_\lambda)\to\Hom(C,\bigoplus_\lambda K_\lambda)\) is bijective. Besides the uniform bound \(\Ext^q(C,N)=0\) for \(q>d_C\) and all constructible \(N\) (Lemma \hyperref[bp20]{5.6} below), we use that \(\Ext^q(C,-)\) commutes with filtered colimits of complexes that vanish below a fixed degree, in particular of sheaves (Lemma \bpout{3.22}, proved in Section \bpout{3.6}).

% object: subsec-A-uniform-bound-for-Ext (blueprint/chap-Construction-of-f/sec-Compactness-of-constructible-complexes/subsec-A-uniform-bound-for-Ext/_section.md), content hidden
\subsubsection*{\texorpdfstring{5.2.1\quad A uniform bound for Ext}{5.2.1 A uniform bound for Ext}}
\phantomsection\label{bp13}%
\addcontentsline{toc}{subsubsection}{\texorpdfstring{5.2.1\quad A uniform bound for Ext}{5.2.1 A uniform bound for Ext}}
% text: prose-05-07
\phantomsection\label{bp14}%
The idea of the bound is to count cohomological dimensions modulo \(\pi_E^n\), add one degree for the derived limit over \(n\), and invert \(\pi_E\) by passing to a colimit. For torsion coefficients the count is direct. By (3.1), \(R\!\operatorname{Hom}(a_!L,N)\) is the cohomology on the stratum \(S\) of a complex obtained from \(N\) by restriction and \(Rj_*\). Each of these steps raises degrees by a bounded amount that does not depend on \(N\). With coefficients \(E\) we cannot count in this way, because \(E\)-sheaves are not torsion and \(\operatorname{cd}_\ell\) only controls torsion sheaves. This is why we choose lattices on the strata. Inverting \(\pi_E\) at the end then needs the continuity of Ext from Section \bpout{3.6}.

The bound \(d_C\) on \(\Ext^q(C,-)\) is a sum of the cohomological dimensions \(\operatorname{cd}_\ell\) of Definition \bpout{2.45}. In the proof we compute \(R\!\operatorname{Hom}(a_!L_0,K)\), for a derived complete \(K\in D(X_{\mathrm{pro}\acute{e}t},\mathcal O_E)\), as a derived inverse limit over the reductions modulo \(\pi_E^n\). For this, the functors in (3.1) have to commute with derived inverse limits, and we use the next lemma with \(\Lambda=\mathcal O_E\).

% object: lemma-limits-sections-support (blueprint/chap-Construction-of-f/sec-Compactness-of-constructible-complexes/subsec-A-uniform-bound-for-Ext/lemma-limits-sections-support/_section.md)
\phantomsection\label{bp15}%
\begin{bpstatementit}{Lemma~5.4}{ (Limits and sections with support)}
Let \(\Lambda\) be admissible, let \(i:Z\hookrightarrow X\) be a closed immersion with open complement \(j:V\hookrightarrow X\), and let \(u:S\hookrightarrow Z\) be an open immersion. The functors \(R\underline\Gamma_Z:D(X_{\mathrm{pro}\acute{e}t},\Lambda)\to D(Z_{\mathrm{pro}\acute{e}t},\Lambda)\) and \(u^*R\underline\Gamma_Z:D(X_{\mathrm{pro}\acute{e}t},\Lambda)\to D(S_{\mathrm{pro}\acute{e}t},\Lambda)\) are left \(t\)-exact and commute with derived inverse limits of sequences, and so does tensoring with a finite locally free \(\Lambda_S\)-module.
\end{bpstatementit}

% proof: proof/lemma-limits-sections-support
\begin{proof}[{Proof}]\phantomsection\label{bp16}%
The functor \(R\underline\Gamma_Z\) is left \(t\)-exact because it is right adjoint to the exact functor \(i_*\). Restricting along the open immersion \(u\) is exact, so \(u^*R\underline\Gamma_Z\) is left \(t\)-exact as well.

For the limits, recall that the derived inverse limit of a sequence \((K_n)\) is the fibre of the map \(\prod_nK_n\to\prod_nK_n\) given by the identity minus the transition maps. It therefore suffices to show that the functors preserve products and fibres. Restriction to an open is exact and preserves limits, so this holds for \(j^*\) and \(u^*\). It holds for \(i^*\) by [Stacks, \href{https://stacks.math.columbia.edu/tag/09AA}{09AA}], and for \(Rj_*\) because \(Rj_*\) is a right adjoint. It follows for \(R\underline\Gamma_Z=i^*\fib(\mathrm{id}\to Rj_*j^*)\). Finally, tensoring with a finite locally free sheaf \(L\) is exact and is right adjoint to tensoring with \(L^\vee\), so it preserves products as well.
\end{proof}

% text: prose-05-08
\phantomsection\label{bp17}%
The main case is that of lattices. There we count cohomological dimensions for the reductions modulo \(\pi_E^n\), and we pass to the limit by Lemma \hyperref[bp15]{5.4}.

% object: lemma-ext-bound-lattices (blueprint/chap-Construction-of-f/sec-Compactness-of-constructible-complexes/subsec-A-uniform-bound-for-Ext/lemma-ext-bound-lattices/_section.md)
\phantomsection\label{bp18}%
\begin{bpstatementit}{Lemma~5.5}{ (Ext bound for lattices)}
Let \(i:Z\hookrightarrow X\) be a closed immersion with open complement \(j:X\setminus Z\hookrightarrow X\), let \(u:S\hookrightarrow Z\) be an open immersion, put \(a=i\circ u\), and let \(L_0\) be a lisse \(\mathcal O_{E,S}\)-module that is locally free of finite rank. Then for every locally closed immersion \(b:T\hookrightarrow X\) and every lisse \(\mathcal O_{E,T}\)-module \(M_0\) that is locally free of finite rank,

\[
\Ext^q_{\mathcal O_E}(a_!L_0,b_!M_0)=0
\quad(q>\operatorname{cd}_\ell(j)+\operatorname{cd}_\ell(S)+2).
\tag{5.1}
\]
\end{bpstatementit}

% proof: proof/lemma-ext-bound-lattices
\begin{proof}[{Proof}]\phantomsection\label{bp19}%
We first bound the complex in (3.1) for the torsion reductions of \(b_!M_0\) by cohomological dimensions, and then pass to the derived limit, which costs one degree. Both cohomological dimensions are finite by [Stacks, \href{https://stacks.math.columbia.edu/tag/0F10}{0F10}] and Lemma \bpout{2.46}.

By Lemma lemma-integral-models and (4), \(b_!M_0\) lies in \(\Dbc(X,\mathcal O_E)\), so it is derived complete. Since \(b_!\) is exact and \(M_0\) is \(\pi_E\)-torsion free, its reductions are

\[
b_!M_0\otimes^{\mathbf L}_{\mathcal O_{E,X}}\mathcal O_{E,X}/\pi_E^n
=b_!(M_0/\pi_E^n).
\]

These are constructible torsion sheaves. They are pulled back from \(X_{\acute{e}t}\), because extension by zero and closed direct image commute with this pullback [Stacks, \href{https://stacks.math.columbia.edu/tag/09AF}{09AF}, \href{https://stacks.math.columbia.edu/tag/09A6}{09A6}].

\emph{Torsion targets.} We bound the complex that appears in (3.1) when the target is one of the torsion sheaves \(b_!(M_0/\pi_E^n)\). By [Stacks, \href{https://stacks.math.columbia.edu/tag/09A6}{09A6}], \(Rj_*\) in \(D(X_{\mathrm{pro}\acute{e}t},\mathcal O_E/\pi_E^n)\) of a sheaf pulled back from the étale site is the pullback of its direct image in \(D(X_{\acute{e}t},\mathcal O_E/\pi_E^n)\). On the étale site the degrees are controlled by the cohomological dimension, and hence \(Rj_*j^*b_!(M_0/\pi_E^n)\) has cohomology in degrees \([0,\operatorname{cd}_\ell(j)]\). The fibre of the map from a sheaf to this complex lives in degrees \([0,\operatorname{cd}_\ell(j)+1]\), so the fibre formula of Lemma lemma-sections-support-closed-subscheme shows that

\[
R\underline\Gamma_Z\,b_!(M_0/\pi_E^n)\in D^{[0,\operatorname{cd}_\ell(j)+1]}(Z_{\mathrm{pro}\acute{e}t},\mathcal O_E/\pi_E^n),
\]

and its cohomology sheaves are torsion sheaves pulled back from the étale site, though not necessarily constructible. Restricting to \(S\) and tensoring with the locally free lisse sheaf \(L_0^\vee/\pi_E^n\) does not change this range of degrees. Since the cohomology of \(S\) with coefficients in a torsion sheaf vanishes above \(\operatorname{cd}_\ell(S)\), the hypercohomology spectral sequence on \(S\) gives

\[
H^q R\Gamma\bigl(S,(L_0^\vee/\pi_E^n)\otimes u^*R\underline\Gamma_Z\,b_!(M_0/\pi_E^n)\bigr)=0
\quad(q>\operatorname{cd}_\ell(j)+\operatorname{cd}_\ell(S)+1),
\tag{5.2}
\]

for every \(n\), \(b\) and \(M_0\).

\emph{Integral targets.} For the lattice \(b_!M_0\) itself we write the relevant complex as a derived limit of its reductions modulo \(\pi_E^n\). Consider the complex \(L_0^\vee\otimes u^*R\underline\Gamma_Z\,b_!M_0\). By Lemma \hyperref[bp15]{5.4}, the functors that form it commute with derived inverse limits of sequences. They also commute with reduction modulo \(\pi_E^n\), which is the cone of multiplication by \(\pi_E^n\). Since \(b_!M_0\) is derived complete, so is this complex, and its reductions are the complexes in (5.2). In other words, the complex is the derived limit of its reductions, and we may compute its cohomology on \(S\) from theirs. Derived global sections preserve products and cones, and hence derived inverse limits. Using (3.1) with \(\mathcal O_E\) coefficients, that is, for \(R\!\operatorname{Hom}\) in \(D(X_{\mathrm{pro}\acute{e}t},\mathcal O_E)\), we get

\[
R\!\operatorname{Hom}_X(a_!L_0,b_!M_0)
=R\Gamma\bigl(S,L_0^\vee\otimes u^*R\underline\Gamma_Z\,b_!M_0\bigr)
=R\varprojlim_n R\Gamma\bigl(S,(L_0^\vee\otimes u^*R\underline\Gamma_Z\,b_!M_0)\otimes^{\mathbf L}\mathcal O_E/\pi_E^n\bigr).
\]

A derived inverse limit of a sequence raises an upper bound on cohomology by at most one. Indeed, it is the fibre of the difference of the identity and the transition maps on the product, and products of abelian groups are exact. This is the Milnor sequence [Stacks, \href{https://stacks.math.columbia.edu/tag/08TB}{08TB}]. Applied to the bound (5.2) for the reductions, it gives (5.1).
\end{proof}

% object: lemma-uniform-ext-bound (blueprint/chap-Construction-of-f/sec-Compactness-of-constructible-complexes/subsec-A-uniform-bound-for-Ext/lemma-uniform-ext-bound/_section.md)
\phantomsection\label{bp20}%
\begin{bpstatementit}{Lemma~5.6}{ (Uniform Ext bound)}
For every \(C\in\Cons(X,\Lambda)\) there is an integer \(d_C\geq0\) such that

\[
\Ext^q(C,N)=0
\qquad\text{for all }q>d_C\text{ and all }N\in\Cons(X,\Lambda).
\]
\end{bpstatementit}

% text: prose-05-09
\phantomsection\label{bp21}%
Suppose \(C=a_!L\), with \(a\) factored as \(S\hookrightarrow Z\hookrightarrow X\) and \(j\) the open complement of \(Z\), and suppose that \(L\) contains a lisse \(\mathcal O_E\)-lattice. By Lemma \bpout{3.24}, this holds after refining \(S\). In this case the proof gives \(d_C=\operatorname{cd}_\ell(j)+\operatorname{cd}_\ell(S)+2\). One of the two extra degrees comes from the cone defining \(R\underline\Gamma_Z\), and the other from a derived inverse limit over the reductions modulo \(\pi_E^n\).

% proof: proof/lemma-uniform-ext-bound
\begin{proof}[{Proof of Lemma 5.6}]\phantomsection\label{bp22}%
We reduce to extensions by zero of lisse sheaves on strata, prove the bound for \(\Lambda=E\) by choosing lattices, applying Lemma \hyperref[bp18]{5.5} and inverting \(\pi_E\), and reduce \(\Lambda=\overline{\mathbf Q}_\ell\) to finite extensions of \(E\).

\emph{Reduction to strata.} Bounds pass to extensions in either variable. Indeed, if \(0\to C'\to C\to C''\to0\) is exact, then in the long exact sequence \(\Ext^q(C'',N)\to\Ext^q(C,N)\to\Ext^q(C',N)\) the outer groups vanish above the larger of the two bounds for \(C'\) and \(C''\), and so the middle one does too. Hence a bound that holds for every subquotient of a finite filtration of the source holds for the source, if we take the maximum. The same sequence in the second variable shows that a bound that holds for every subquotient of a finite filtration of the target holds for the target, with the same value. In view of Lemma lemma-constructible-sheaves, we may therefore assume \(C=a_!L\), with \(a=i\circ u:S\hookrightarrow Z\hookrightarrow X\) as in Lemma \bpout{3.21} and \(L\) lisse. It suffices to find an integer \(d\), depending only on \(a\) and \(L\), with \(\Ext^q(a_!L,b_!M)=0\) for \(q>d\), for all locally closed \(b:T\hookrightarrow X\) and lisse \(M\) on \(T\). Let \(j:X\setminus Z\hookrightarrow X\) be the open complement of \(Z\). When \(L\) contains a lisse \(\mathcal O_E\)-lattice, we show that \(d=\operatorname{cd}_\ell(j)+\operatorname{cd}_\ell(S)+2\) works.

\emph{Lattices.} First take \(\Lambda=E\). Since torsion arguments need integral structures, we choose lattices. After refining the stratifications (Lemma \bpout{3.24}), \(L\) and \(M\) contain lisse \(\mathcal O_E\)-lattices \(L_0\) and \(M_0\). Refining is harmless, because it replaces \(a_!L\) and \(b_!M\) by finite extensions of sheaves of the same form, and the first step handles finite extensions. Moreover, the refinement of the source stratum depends only on \(L\), and that of the target only on \(M\). So the bound obtained for the refined source strata is still independent of \(N\).

The lattices are locally free of finite rank over \(\mathcal O_E\). Indeed, being lisse, they are locally of the form \(\mathcal O_{E,U}\otimes_{\mathcal O_E}P\) with \(P\) finitely generated (Definition \bpout{3.2}). As subsheaves of \(E\)-modules they are torsion free, so \(P\) is free. Lemma \hyperref[bp18]{5.5} now gives \(\Ext^q_{\mathcal O_E}(a_!L_0,b_!M_0)=0\) for \(q>\operatorname{cd}_\ell(j)+\operatorname{cd}_\ell(S)+2\).

\emph{Inverting \(\pi_E\).} Inverting \(\pi_E\) is a filtered colimit. By continuity of Ext out of \(a_!L_0\), the Ext groups over \(E\) are therefore colimits of groups that all vanish in the same range. In detail, we have \(a_!L=a_!L_0[1/\pi_E]\) and \(b_!M=(b_!M_0)[1/\pi_E]=\varinjlim(b_!M_0\xrightarrow{\pi_E}b_!M_0\xrightarrow{\pi_E}\cdots)\). The adjunction between extension and restriction of scalars gives

\[
R\!\operatorname{Hom}_E\bigl(a_!L,b_!M\bigr)
=R\!\operatorname{Hom}_{\mathcal O_E}\bigl(a_!L_0,(b_!M_0)[1/\pi_E]\bigr).
\]

By Lemma lemma-continuity-ext-out-constructible-complex the cohomology of the right side is the colimit of the groups \(\Ext^q_{\mathcal O_E}(a_!L_0,b_!M_0)\). Hence (5.1) implies the bound for \(E\).

\emph{The case \(\Lambda=\overline{\mathbf Q}_\ell\).} The bound for \(E\) depends only on the strata and on \(\operatorname{cd}_\ell\), not on \(E\), so it passes to the colimit over finite extensions. We first fix everything that depends on \(L\) alone, before the target is given. By Lemma \bpout{3.17}, applied to \(L\) alone, there are a finite extension \(E_0/\mathbf Q_\ell\) and a lisse \(E_0\)-sheaf \(L'\) with \(L=L'_{\overline{\mathbf Q}_\ell}\). After refining \(S\) (Lemma \bpout{3.24}), \(L'\) contains a lisse \(\mathcal O_{E_0}\)-lattice \(L'_0\). As in the case \(\Lambda=E\), the refinement replaces \(a_!L\) by a finite extension of sheaves of the same form, and it depends only on \(L\). For every finite extension \(E'\supset E_0\), the sheaf \(L'_0\otimes_{\mathcal O_{E_0}}\mathcal O_{E'}\) is a lisse \(\mathcal O_{E'}\)-lattice in \((L')_{E'}\) on the same strata. The case \(\Lambda=E'\) therefore gives \(\Ext^q_{E'}\bigl((a_!L')_{E'},N'\bigr)=0\) for \(q>\operatorname{cd}_\ell(j)+\operatorname{cd}_\ell(S)+2\), for each refined stratum \(S\) and every \(N'\in\Cons(X,E')\). This bound does not depend on \(E'\) or on \(N'\).

Now let \(b\) and \(M\) be given. By Lemma \bpout{3.17} again, applied to \(M\), there are a finite extension \(E\supset E_0\) (enlarge the field given by the lemma to contain \(E_0\)) and a lisse \(E\)-sheaf \(M'\) with \(M=M'_{\overline{\mathbf Q}_\ell}\). The field \(E\) depends on \(M\), but neither the strata of the source nor the lattice \(L'_0\) do. Then \(a_!L=(a_!L')_{\overline{\mathbf Q}_\ell}\) and \(b_!M=(b_!M')_{\overline{\mathbf Q}_\ell}\), and by Lemma lemma-descent-finite-coefficient-field,

\[
\Ext^q_{\overline{\mathbf Q}_\ell}(a_!L,b_!M)
=\varinjlim_{E'\supset E}\Ext^q_{E'}\bigl((a_!L')_{E'},(b_!M')_{E'}\bigr).
\]

Every group on the right vanishes for \(q>\operatorname{cd}_\ell(j)+\operatorname{cd}_\ell(S)+2\) by the previous paragraph, and hence so does the colimit.
\end{proof}

% object: example-ext-bound-finite-field (blueprint/chap-Construction-of-f/sec-Compactness-of-constructible-complexes/subsec-A-uniform-bound-for-Ext/example-ext-bound-finite-field/_section.md)
\phantomsection\label{bp23}%
\begin{bpstatement}{Example~5.7}{ (The bound over a finite field)}
The proof gives an explicit bound, which is far from sharp. For instance, for \(X=\operatorname{Spec}\mathbf F_q\) and \(C=N=E\) we have \(Z=S=X\). The complement of \(Z\) is then empty, so \(R\underline\Gamma_Z\) is the identity. Moreover \(\operatorname{cd}_\ell(\mathbf F_q)=1\), and so the formula, read with \(\operatorname{cd}_\ell(j)=0\) for the empty open, gives \(d_C=0+1+2=3\), while \(\Ext^q(E,E)\) vanishes already for \(q>1\) (Example \hyperref[bp11]{5.3}). The two wasted degrees do no harm, because the arguments below only need some bound independent of \(N\).
\end{bpstatement}


% object: subsec-Ind-constructible-sheaves (blueprint/chap-Construction-of-f/sec-Compactness-of-constructible-complexes/subsec-Ind-constructible-sheaves/_section.md), content hidden
\subsubsection*{\texorpdfstring{5.2.2\quad Ind-constructible sheaves}{5.2.2 Ind-constructible sheaves}}
\phantomsection\label{bp24}%
\addcontentsline{toc}{subsubsection}{\texorpdfstring{5.2.2\quad Ind-constructible sheaves}{5.2.2 Ind-constructible sheaves}}
% text: prose-05-10
\phantomsection\label{bp25}%
The targets in the compactness statement are direct sums of objects of \(\Loc(\Dbc(X,\Lambda))\), and their cohomology sheaves need not be constructible. Already the sheaf \(\bigoplus_{n\in\mathbf N}\Lambda_X\) is not constructible, since its stalks are not finitely generated. Such sheaves are, however, filtered colimits of constructible sheaves, and both continuity and the uniform bound for Ext groups extend to these colimits.

Recall that the ind-category \(\operatorname{Ind}(\mathcal C)\) of a small category \(\mathcal C\) has as objects the filtered diagrams in \(\mathcal C\) (“formal filtered colimits”), with \(\Hom(\varinjlim_\alpha C_\alpha,\varinjlim_\beta D_\beta)
=\varprojlim_\alpha\varinjlim_\beta\Hom(C_\alpha,D_\beta)\). Proposition \hyperref[bp27]{5.9} below says that for constructible sheaves the formal and the actual colimits agree.

The dévissage arguments need a class of sheaves that is abelian and closed under extensions and direct sums, and on which \(\Ext^q(C,-)\) behaves as on constructible sheaves. By Proposition \hyperref[bp27]{5.9}, the category \(\IndCons(X,\Lambda)\) has the first three properties, and Corollary \hyperref[bp30]{5.10} gives the last.

% object: definition-ind-constructible-sheaves (blueprint/chap-Construction-of-f/sec-Compactness-of-constructible-complexes/subsec-Ind-constructible-sheaves/definition-ind-constructible-sheaves/_section.md)
\phantomsection\label{bp26}%
\begin{bpstatement}{Definition~5.8}{ (Ind-constructible sheaves)}
An \emph{ind-constructible sheaf} on \(X\) is a sheaf of \(\Lambda_X\)-modules on \(X_{\mathrm{pro}\acute{e}t}\) that is a filtered colimit, computed in \(\operatorname{Mod}(X_{\mathrm{pro}\acute{e}t},\Lambda)\), of constructible sheaves. We write \(\IndCons(X,\Lambda)\subset\operatorname{Mod}(X_{\mathrm{pro}\acute{e}t},\Lambda)\) for the full subcategory of ind-constructible sheaves.
\end{bpstatement}

% object: proposition-ind-constructible-sheaves (blueprint/chap-Construction-of-f/sec-Compactness-of-constructible-complexes/subsec-Ind-constructible-sheaves/proposition-ind-constructible-sheaves/_section.md)
\phantomsection\label{bp27}%
\begin{bpstatementit}{Proposition~5.9}{ (Ind-constructible sheaves)}
\begin{enumerate}
\item The functor \(\operatorname{Ind}(\Cons(X,\Lambda))\to\operatorname{Mod}(X_{\mathrm{pro}\acute{e}t},\Lambda)\) sending a formal filtered colimit to the actual colimit is fully faithful, with essential image \(\IndCons(X,\Lambda)\).
\item \(\IndCons(X,\Lambda)\) is closed in \(\operatorname{Mod}(X_{\mathrm{pro}\acute{e}t},\Lambda)\) under filtered colimits, direct sums, kernels, cokernels and extensions. In particular it is a Grothendieck abelian category and its inclusion is exact.
\end{enumerate}
\end{bpstatementit}

% proof: proof/proposition-ind-constructible-sheaves
\begin{proof}[{Proof}]\phantomsection\label{bp28}%
The proof uses one property of constructible sheaves: maps out of them commute with filtered colimits. Because of this, the formal colimits of the ind-category behave like actual colimits. Each closure property in (2) then reduces to the same property of \(\Cons(X,\Lambda)\), once the objects and maps involved are written as filtered colimits of constructible sheaves and maps between them.

\emph{Finite presentation.} Constructible sheaves are finitely presented in \(\operatorname{Mod}(X_{\mathrm{pro}\acute{e}t},\Lambda)\). Indeed, Lemma \bpout{3.22} in degree zero shows that \(\Hom(C,-)\) commutes with filtered colimits for \(C\) constructible. By Lemma lemma-constructible-sheaves,(2), the category \(\Cons(X,\Lambda)\) is essentially small, abelian, and closed under extensions. We fix a small skeleton, so that the ind-category in (1) is defined.

\emph{Full faithfulness.} For filtered diagrams \((C_\alpha)\) and \((D_\beta)\) of constructible sheaves, finite presentation implies

\[
\Hom(\varinjlim_\alpha C_\alpha,\varinjlim_\beta D_\beta)
=\varprojlim_\alpha\varinjlim_\beta\Hom_{\Cons(X,\Lambda)}(C_\alpha,D_\beta),
\]

which is the Hom set in the ind-category. This proves (1), since the essential image is \(\IndCons(X,\Lambda)\) by definition.

\emph{Canonical presentations.} The closure properties below involve several ind-constructible sheaves at once, so to compare their presentations we use a presentation that involves no choices. Namely, for \(B\in\IndCons(X,\Lambda)\), let \(J(B)\) be the category of constructible sheaves \(C\) with a map \(C\to B\); we show that \(B=\varinjlim_{J(B)}C\).

The category \(J(B)\) is filtered. Indeed, two objects \(C_1\to B\) and \(C_2\to B\) both map to \(C_1\oplus C_2\to B\), and two parallel arrows \(C\rightrightarrows C'\) over \(B\) are equalized by the cokernel of their difference, which is constructible and still maps to \(B\). Now write \(B=\varinjlim_\alpha C_\alpha\). By finite presentation, every \(C\to B\) factors through some \(C_\alpha\), and two factorizations agree at a later stage. This means that the given diagram is cofinal in \(J(B)\), and hence \(B=\varinjlim_{J(B)}C\).

\emph{Filtered colimits and direct sums.} Let \(B=\varinjlim_\lambda B_\lambda\) be a filtered colimit with \(B_\lambda\in\IndCons(X,\Lambda)\). To present \(B\), consider the category of pairs \((\lambda,C\to B_\lambda)\) with \(C\) constructible. It is filtered: we move finitely many indices to a common upper bound, and then use direct sums and cokernels in \(J(B_\lambda)\). Moreover, the colimit of the \(C\) over this category is \(\varinjlim_\lambda\varinjlim_{J(B_\lambda)}C=B\). A direct sum is the filtered colimit of its finite subsums, so it lies in \(\IndCons(X,\Lambda)\) as well.

\emph{Kernels and cokernels.} Let \(f:B\to B'\) be a map in \(\IndCons(X,\Lambda)\). It suffices to write \(f\) as a filtered colimit of maps \(v\) of constructible sheaves. Since filtered colimits are exact, \(\ker f\) and \(\coker f\) are then the filtered colimits of the constructible sheaves \(\ker v\) and \(\coker v\). Let \(P_f\) be the category of commutative squares

\[
\begin{matrix} C&\xrightarrow{v}&D\\
\downarrow&&\downarrow\\ B&\xrightarrow{f}&B',\end{matrix}
\qquad C,D\in\Cons(X,\Lambda).
\]

This category is filtered, because finite colimits of such squares are formed by finite colimits of the arrows \(v\) in the abelian category \(\Cons(X,\Lambda)\). We claim that the functor \(P_f\to J(B)\), \((v)\mapsto(C\to B)\), is cofinal. Indeed, every \(C\to B\) is the source of a square, namely \(v=\mathrm{id}_C\) with \(C\to B'\) the composite. If \(C\to C_1\xrightarrow{v_1}D_1\) and \(C\to C_2\xrightarrow{v_2}D_2\) are two squares under the same \(C\to B\), then both map to the square with source \(C_1\amalg_CC_2\) and target \(D_1\amalg_CD_2\). Here the pushouts are taken in \(\Cons(X,\Lambda)\), and they map to \(B\) and \(B'\) because the two maps from \(C\) to \(B\), and the two maps from \(C\) to \(B'\), agree.

Similarly, the functor \(P_f\to J(B')\), \((v)\mapsto(D\to B')\), is cofinal. Every \(D\to B'\) is the target of a square, namely \(v:0\to D\). If \(C_1\xrightarrow{v_1}D_1\) and \(C_2\xrightarrow{v_2}D_2\) are two squares under the same \(D\to B'\), then both map to the square with source \(C_1\oplus C_2\) and target \(D_1\amalg_DD_2\). The two cofinality statements give \(f=\varinjlim_{P_f}v\), and exactness of filtered colimits gives

\[
\ker f=\varinjlim_{P_f}\ker v,\qquad
\coker f=\varinjlim_{P_f}\coker v,
\]

both of which lie in \(\IndCons(X,\Lambda)\).

\emph{Extensions.} Let \(0\to N\to M\to B\to0\) be exact in \(\operatorname{Mod}(X_{\mathrm{pro}\acute{e}t},\Lambda)\), with \(N=\varinjlim_\beta D_\beta\) and \(B=\varinjlim_\alpha C_\alpha\) ind-constructible. We show that \(M\) is a filtered colimit of extensions of constructible sheaves. First, the pullback of the extension along \(C_\alpha\to B\) is ind-constructible. Its class lies in \(\Ext^1(C_\alpha,N)\), and Lemma \bpout{3.22} gives

\[
\Ext^1(C_\alpha,N)
=\varinjlim_\beta\Ext^1(C_\alpha,D_\beta).
\]

So the pulled-back extension is pushed out from an extension \(M_{\alpha\beta}\) of \(C_\alpha\) by some \(D_\beta\). Since \(\Cons(X,\Lambda)\) is closed under extensions, \(M_{\alpha\beta}\) and each of its pushouts along \(D_\beta\to D_{\beta'}\) are constructible. Their colimit over \(\beta'\) is the pullback of \(M\) to \(C_\alpha\). Second, since filtered colimits are exact, \(M\) is the colimit over \(\alpha\) of these pullbacks. By closure under filtered colimits, \(M\in\IndCons(X,\Lambda)\).

\emph{Grothendieck.} The category \(\IndCons(X,\Lambda)\) is Grothendieck abelian in the sense of [Stacks, \href{https://stacks.math.columbia.edu/tag/079A}{079A}]. Indeed, by the above its colimits exist and are computed in \(\operatorname{Mod}(X_{\mathrm{pro}\acute{e}t},\Lambda)\), so filtered colimits are exact. The direct sum of the objects of the skeleton of \(\Cons(X,\Lambda)\) is a generator.
\end{proof}

% text: prose-05-11
\phantomsection\label{bp29}%
Since every ind-constructible sheaf is a filtered colimit of constructible ones, the continuity of \(\Ext^q(C,-)\) lets us extend the uniform bound of Lemma \hyperref[bp20]{5.6} to the whole category \(\IndCons(X,\Lambda)\). Writing a direct sum as the filtered colimit of its finite subsums, we also see that \(\Ext^q(C,-)\) commutes with direct sums of ind-constructible sheaves.

% object: corollary-ind-heart-b (blueprint/chap-Construction-of-f/sec-Compactness-of-constructible-complexes/subsec-Ind-constructible-sheaves/corollary-ind-heart-b/_section.md)
\phantomsection\label{bp30}%
\begin{bpstatementit}{Corollary~5.10}{ (Ext from a constructible sheaf into ind-constructible sheaves)}
Let \(C\in\Cons(X,\Lambda)\), with \(d_C\) as in Lemma \hyperref[bp20]{5.6}. Then \(\Ext^q(C,B)=0\) for all \(q>d_C\) and all \(B\in\IndCons(X,\Lambda)\), and each functor \(\Ext^q(C,-)\) commutes with direct sums of objects of \(\IndCons(X,\Lambda)\).
\end{bpstatementit}

% proof: proof/corollary-ind-heart-b
\begin{proof}[{Proof}]\phantomsection\label{bp31}%
Write \(B=\varinjlim_\alpha B_\alpha\) as a filtered colimit of constructible sheaves, which is possible by Proposition \hyperref[bp27]{5.9}. Sheaves are complexes concentrated in degree \(0\), so they vanish below a common degree, and Lemma \bpout{3.22} applies. It shows that \(\Ext^q(C,B)=\varinjlim_\alpha\Ext^q(C,B_\alpha)\). Each group on the right vanishes for \(q>d_C\) by Lemma \hyperref[bp20]{5.6}, and so does the colimit.

For the second claim, let \((B_\lambda)\) be a family in \(\IndCons(X,\Lambda)\). Its direct sum is the filtered colimit of the finite subsums, and \(\Ext^q(C,-)\) commutes with finite direct sums because it is additive. Applying Lemma \bpout{3.22} once more, to the diagram of finite subsums, we obtain \(\Ext^q(C,\bigoplus_\lambda B_\lambda)=\bigoplus_\lambda\Ext^q(C,B_\lambda)\).
\end{proof}

% text: prose-05-12
\phantomsection\label{bp32}%
Objects of \(\Loc(\Dbc(X,\Lambda))\) can be unbounded, and we do not know a priori that their truncations stay in \(\Loc(\Dbc(X,\Lambda))\). For this reason, here and in Section \hyperref[bp53]{5.4}, we carry out truncation arguments among complexes with ind-constructible cohomology sheaves.

% object: corollary-ambient-ind-derived (blueprint/chap-Construction-of-f/sec-Compactness-of-constructible-complexes/subsec-Ind-constructible-sheaves/corollary-ambient-ind-derived/_section.md)
\phantomsection\label{bp33}%
\begin{bpstatementit}{Corollary~5.11}{ (The derived category of ind-constructible sheaves)}
The full subcategory

\[
\Dic(X,\Lambda)
=\{K\in D(X_{\mathrm{pro}\acute{e}t},\Lambda):H^n(K)\in\IndCons(X,\Lambda)\text{ for every }n\}
\]

of \(D(X_{\mathrm{pro}\acute{e}t},\Lambda)\) is triangulated, closed under direct sums, and stable under the truncations \(\tau_{\leq n}\), \(\tau_{\geq n}\). In particular \(\Loc(\Dbc(X,\Lambda))\subset\Dic(X,\Lambda)\).
\end{bpstatementit}

% proof: proof/corollary-ambient-ind-derived
\begin{proof}[{Proof}]\phantomsection\label{bp34}%
In the long exact cohomology sequence of a triangle, each cohomology sheaf of the third term is an extension of a kernel by a cokernel of maps between cohomology sheaves of the other two. Since \(\IndCons(X,\Lambda)\) is closed under kernels, cokernels and extensions (Proposition \hyperref[bp27]{5.9}), the third term lies in \(\Dic(X,\Lambda)\) whenever the other two do. Closure under shifts is clear. As direct sums of sheaves are exact, \(H^n\) commutes with direct sums. Since \(\IndCons(X,\Lambda)\) is closed under direct sums, it follows that \(\Dic(X,\Lambda)\) is closed under direct sums. Truncation only removes cohomology sheaves and keeps the others. Hence \(\Dic(X,\Lambda)\) is a localizing subcategory of \(D(X_{\mathrm{pro}\acute{e}t},\Lambda)\) containing \(\Dbc(X,\Lambda)\). It therefore contains the smallest one, \(\Loc(\Dbc(X,\Lambda))\).
\end{proof}


% object: subsec-Rational-coefficients (blueprint/chap-Construction-of-f/sec-Compactness-of-constructible-complexes/subsec-Rational-coefficients/_section.md), content hidden
\subsubsection*{\texorpdfstring{5.2.3\quad Rational coefficients}{5.2.3 Rational coefficients}}
\phantomsection\label{bp35}%
\addcontentsline{toc}{subsubsection}{\texorpdfstring{5.2.3\quad Rational coefficients}{5.2.3 Rational coefficients}}
% text: prose-05-13
\phantomsection\label{bp36}%
On ind-constructible sheaves, \(\Ext^q(C,-)\) vanishes above \(d_C\) and commutes with direct sums (Corollary \hyperref[bp30]{5.10}). Because the bound \(d_C\) is uniform, maps from a constructible sheaf \(C\) into an unbounded complex depend only on its cohomology sheaves in degrees \([-d_C,0]\). This is the step that fails over \(\mathbf R\) in the example at the beginning of this chapter, where \(\Ext^n(\mathbf F_2,\mathbf F_2)\neq0\) for all \(n\). Since the window has the same length for a whole family of targets, compactness follows from the case of sheaves by induction on the length of the window.

% object: lemma-reduction-finitely-many-degrees (blueprint/chap-Construction-of-f/sec-Compactness-of-constructible-complexes/subsec-Rational-coefficients/lemma-reduction-finitely-many-degrees/_section.md)
\phantomsection\label{bp37}%
\begin{bpstatementit}{Lemma~5.12}{ (Reduction to finitely many degrees)}
Let \(C\in\Cons(X,\Lambda)\), with \(d_C\) as in Lemma \hyperref[bp20]{5.6}, and let \(K\in\Dic(X,\Lambda)\) and \(m\in\mathbf Z\). The truncation maps induce a natural isomorphism

\[
\Hom(C,K[m])
\simeq\Hom\bigl(C,(\tau_{\geq m-d_C}\tau_{\leq m}K)[m]\bigr).
\tag{5.3}
\]
\end{bpstatementit}

% proof: proof/lemma-reduction-finitely-many-degrees
\begin{proof}[{Proof}]\phantomsection\label{bp38}%
Cutting off the cohomology above degree \(m\) is harmless for degree reasons alone. To remove the cohomology below degree \(m-d_C\) we use the uniform bound: for bounded \(K\) it kills the tail layer by layer, and for \(K\) unbounded below we write \(K\) as the derived limit of its truncations and control the resulting \(\varprojlim^1\) term.

Replacing \(K\) by \(K[m]\), we may assume \(m=0\). The truncation \(\tau_{>0}K\) lies in \(D^{\geq1}\), and \(C\) is a sheaf in degree \(0\), so there are no nonzero maps from \(C\) to an object of \(D^{\geq1}\). Hence \(\Hom(C,\tau_{>0}K[s])=0\) for \(s=0,-1\). In the long exact sequence of \(\Hom(C,-)\) for the truncation triangle \(\tau_{\leq0}K\to K\to\tau_{>0}K\to\), these are the two groups adjacent to \(\Hom(C,K)\), and so \(\Hom(C,K)=\Hom(C,\tau_{\leq0}K)\). To remove the cohomology in degrees \(\leq-d_C-1\) from \(K'=\tau_{\leq0}K\), we use the triangle \(\tau_{\leq-d_C-1}K'\to K'\to\tau_{\geq-d_C}K'\to\). It suffices to show that

\[
\Hom(C,L)=0\qquad
\text{for every }L\in\Dic(X,\Lambda)\text{ with }H^q(L)=0\text{ for }q>-d_C-1,
\tag{5.4}
\]

applied to \(L=\tau_{\leq-d_C-1}K'\) and to \(L=\tau_{\leq-d_C-1}K'[1]\).

We first prove (5.4) for bounded \(L\). Such an \(L\) is a finite extension of its shifted cohomology sheaves \(H^q(L)[-q]\), and by the long exact Hom sequences it suffices to show that each of them receives no nonzero map from \(C\). A layer \(H^q(L)[-q]\) with \(q\leq-d_C-1\) contributes \(\Hom(C,H^q(L)[-q])=\Ext^{-q}(C,H^qL)\). Since \(-q>d_C\) and \(H^q(L)\) is ind-constructible, this group vanishes by Corollary \hyperref[bp30]{5.10}. For \(L\) unbounded below we cannot argue in finitely many steps, and we use instead that the derived category is left complete.

The pro-étale site has a basis of weakly contractible objects (Theorem \bpout{2.17}). For this reason Proposition \bpout{2.23} applies: products of \(\Lambda_X\)-modules are exact, products in \(D(X_{\mathrm{pro}\acute{e}t},\Lambda)\) are computed termwise, and

\[
L=R\varprojlim_n\tau_{\geq-n}L
\]

for every \(L\). Concretely, if \(L^\bullet\) is a complex representing \(L\), the limit is the derived inverse limit of the tower of good truncations \(\tau_{\geq-n}L^\bullet\) with their canonical transition maps. Since \(\Hom(C,-)\) takes products to products, the triangle defining \(R\varprojlim\) yields the Milnor exact sequence

\[
0\longrightarrow\varprojlim_n{}^1
\Hom(C,\tau_{\geq-n}L[-1])
\longrightarrow\Hom(C,L)
\longrightarrow\varprojlim_n\Hom(C,\tau_{\geq-n}L)
\longrightarrow0.
\]

Each \(\tau_{\geq-n}L\) is bounded with cohomology in degrees \(\leq-d_C-1\), and hence the right-hand term vanishes by the bounded case. For the left-hand term, the triangle

\[
H^{-n-1}(L)[n+1]\longrightarrow \tau_{\geq-n-1}L
\longrightarrow \tau_{\geq-n}L\longrightarrow H^{-n-1}(L)[n+2]
\]

gives an exact sequence

\[
\Ext^{n}(C,H^{-n-1}L)\to\Hom(C,\tau_{\geq-n-1}L[-1])\to\Hom(C,\tau_{\geq-n}L[-1])
\to\Ext^{n+1}(C,H^{-n-1}L).
\]

The outer terms vanish for \(n>d_C\) by Corollary \hyperref[bp30]{5.10}. The transition maps are therefore isomorphisms for \(n>d_C\), so the tower is eventually constant and its \(\varprojlim^1\) vanishes. This proves (5.4).
\end{proof}

% text: prose-05-14
\phantomsection\label{bp39}%
The window \([m-d_C,m]\) in (5.3) depends on \(C\) and \(m\) only, not on \(K\), and so we can treat a whole family \((K_\lambda)\) at once in the compactness theorem.

% object: theorem-compactness-rational-coefficients (blueprint/chap-Construction-of-f/sec-Compactness-of-constructible-complexes/subsec-Rational-coefficients/theorem-compactness-rational-coefficients/_section.md)
\phantomsection\label{bp40}%
\begin{bpstatementit}{Theorem~5.13}{ (Compactness with rational coefficients)}
For \(\Lambda=E\) or \(\overline{\mathbf Q}_\ell\), every object of \(\Dbc(X,\Lambda)\) is compact in \(\Loc(\Dbc(X,\Lambda))\). More precisely, for \(C\in\Dbc(X,\Lambda)\) and every family \((K_\lambda)\) of objects of \(\Dic(X,\Lambda)\), the natural map

\[
\bigoplus_\lambda\Hom(C,K_\lambda)
\longrightarrow\Hom\Bigl(C,\bigoplus_\lambda K_\lambda\Bigr)
\]

is an isomorphism.
\end{bpstatementit}

% proof: proof/theorem-compactness-rational-coefficients
\begin{proof}[{Proof}]\phantomsection\label{bp41}%
The idea is that, by Lemma \hyperref[bp37]{5.12}, maps out of a constructible sheaf depend only on the cohomology of the target in \(d_C+1\) consecutive degrees, uniformly in the family. For targets with cohomology in a fixed bounded interval the claim follows from the case of ind-constructible sheaves by induction on the length.

The first assertion follows from the second, since \(\Loc(\Dbc(X,\Lambda))\subset\Dic(X,\Lambda)\) (Corollary \hyperref[bp33]{5.11}) and direct sums in \(\Loc(\Dbc(X,\Lambda))\) are those of \(D(X_{\mathrm{pro}\acute{e}t},\Lambda)\).

First let \(C\in\Cons(X,\Lambda)\), and suppose that all \(K_\lambda\) have cohomology in a fixed interval \([a,b]\). We prove the claim, for all shifts of the \(K_\lambda\) at once, by induction on \(b-a\). If \(a=b\), then \(K_\lambda=B_\lambda[-a]\) with \(B_\lambda\in\IndCons(X,\Lambda)\), and the claim is that \(\Ext^q(C,-)\) commutes with direct sums of objects of \(\IndCons(X,\Lambda)\), which is Corollary \hyperref[bp30]{5.10}. In general, the triangles \(\tau_{\leq b-1}K_\lambda\to K_\lambda\to H^b(K_\lambda)[-b]\to\) are compatible with direct sums, because \(H^n\) commutes with direct sums. We apply \(\Hom(C,-)\) to these triangles and take the direct sum over \(\lambda\); since direct sums of abelian groups are exact, the result is again a long exact sequence. It maps to the long exact sequence of the summed triangle. The outer terms are of the form covered by the induction hypothesis (cohomology in \([a,b-1]\) and in \([b,b]\)), and so the five lemma gives the claim for cohomology in \([a,b]\).

For arbitrary \(K_\lambda\), we apply Lemma \hyperref[bp37]{5.12} with \(m=0\), which allows us to replace each \(K_\lambda\), and \(\bigoplus_\lambda K_\lambda\), by its truncation to degrees \([-d_C,0]\). Since the interval is the same for the whole family and truncation commutes with direct sums, the bounded case applies.

From constructible sheaves \(C\) we pass to complexes by dévissage. For fixed \((K_\lambda)\), the objects \(C\) for which the map is an isomorphism are closed under shifts. They are closed under cones by the five lemma, applied to the long exact sequences of \(\Hom(-,K_\lambda)\) summed over \(\lambda\) and of \(\Hom(-,\bigoplus_\lambda K_\lambda)\). A bounded constructible complex is a finite extension of shifts of its constructible cohomology sheaves, and hence it belongs to this class.
\end{proof}

% text: prose-05-15
\phantomsection\label{bp42}%
Lemma \hyperref[bp97]{5.30} combines this theorem with the corresponding statements for the other coefficient rings and shows that a small skeleton of \(D^b_{\mathrm{ctf}}(X,\Lambda)\) is a set of compact generators of \(\Loc(D^b_{\mathrm{ctf}}(X,\Lambda))\).



% object: sec-Compactness-for-finite-coefficients (blueprint/chap-Construction-of-f/sec-Compactness-for-finite-coefficients/_section.md), content hidden
\subsection*{\texorpdfstring{5.3\quad Compactness for finite coefficients}{5.3 Compactness for finite coefficients}}
\phantomsection\label{bp43}%
\addcontentsline{toc}{subsection}{\texorpdfstring{5.3\quad Compactness for finite coefficients}{5.3 Compactness for finite coefficients}}
% text: prose-05-16
\phantomsection\label{bp44}%
In this section the coefficient ring is a finite admissible ring \(A\). Proper direct image preserves \(D^b_{\mathrm{ctf}}\) by Proposition \bpout{4.35}, and we now show that the constructible generators are compact and that proper direct image preserves direct sums on the localizing categories. In both cases we reduce to the classical étale theory, using bounded flat constructible models and the fact that étale cohomology commutes with direct sums of unbounded complexes.

Finite Tor amplitude is a property of the complex, not of its cohomology sheaves: the cohomology sheaves of an object of \(D^b_{\mathrm{ctf}}\) can have infinite Tor amplitude (Section \bpout{3.4}). The property can, however, be read off from a flat model. This is why we work with the bounded flat models of Lemma \bpout{3.13}.

% object: proposition-compactness-finite-coefficients (blueprint/chap-Construction-of-f/sec-Compactness-for-finite-coefficients/proposition-compactness-finite-coefficients/_section.md)
\phantomsection\label{bp45}%
\begin{bpstatementit}{Proposition~5.14}{ (Compactness for finite coefficients)}
Every object of \(D^b_{\mathrm{ctf}}(X,A)\) is compact in \(\Loc(D^b_{\mathrm{ctf}}(X,A))\), and a small skeleton of \(D^b_{\mathrm{ctf}}(X,A)\) is a set of compact generators. Moreover, \(\Loc(D^b_{\mathrm{ctf}}(X,A))^c=D^b_{\mathrm{ctf}}(X,A)\).
\end{bpstatementit}

% proof: proof/proposition-compactness-finite-coefficients
\begin{proof}[{Proof}]\phantomsection\label{bp46}%
The plan is to prove compactness on the étale site against arbitrary objects of \(D(X_{\acute{e}t},A)\), where the commutation of cohomology with direct sums that we need is classical, and then transfer the result to \(X_{\mathrm{pro}\acute{e}t}\) by \(\nu_X^*\).

Recall that compact objects are closed under shifts, finite extensions and direct summands. Hence, by Lemma \bpout{3.13}, it suffices to treat \(a_!L\). To do so, factor \(a=i\circ u\) with \(i:Z\hookrightarrow X\) closed and \(u:S\hookrightarrow Z\) open, and let \(j:X\setminus Z\hookrightarrow X\) be the complementary open immersion. The idea is then to write \(R\!\operatorname{Hom}(a_!L,-)\) as a composite of functors each of which commutes with direct sums. Since \(L\) has finite projective stalks, it is dualizable with dual \(L^\vee=\underline{\operatorname{Hom}}_A(L,A_S)\), and the étale form of Lemma lemma-sections-support-closed-subscheme shows that for \(K\in D(X_{\acute{e}t},A)\) we have

\[
R\!\operatorname{Hom}_X(a_!L,K)
=R\Gamma\bigl(S,L^\vee\otimes_A u^*R\underline\Gamma_ZK\bigr),
\qquad R\underline\Gamma_ZK=i^*\fib\bigl(K\to Rj_*j^*K\bigr),
\]

where \(R\underline\Gamma_Z\) denotes sections with support in \(Z\). We claim that every functor on the right, viewed as a functor between étale derived categories, preserves arbitrary direct sums. For \(Rj_*\) and \(R\Gamma(S,-)\), this is [Stacks, \href{https://stacks.math.columbia.edu/tag/0F11}{0F11}], which says that a derived direct image of finite cohomological dimension commutes with direct sums of unbounded complexes. It applies to \(A\)-modules because derived direct image and global sections commute with forgetting the coefficient ring [Stacks, \href{https://stacks.math.columbia.edu/tag/0D6J}{0D6J}]. The remaining functors are inverse images, tensor products with finite projective sheaves, and a fibre, and these also preserve direct sums: inverse images are left adjoints, tensor products commute with direct sums, and the fibre of a map between functors that preserve direct sums again preserves them. Hence \(R\!\operatorname{Hom}_X(a_!L,-)\) commutes with direct sums. Taking \(H^0\), we see that \(\Hom(a_!L,-)\) commutes with direct sums, that is, \(a_!L\) is compact. By Lemma \bpout{3.13} it follows that every object of \(D^b_{\mathrm{ctf}}(X_{\acute{e}t},A)\) is compact.

On the pro-étale site, Corollary \bpout{2.44} and Proposition \bpout{2.49} show that the inverse image \(\nu_X^*:D(X_{\acute{e}t},A)\to D(X_{\mathrm{pro}\acute{e}t},A)\) is exact, fully faithful and preserves direct sums. Such a functor identifies the localizing subcategories generated by \(D^b_{\mathrm{ctf}}\) on the two sites. Compactness against all étale complexes therefore implies compactness in \(\Loc(D^b_{\mathrm{ctf}}(X,A))\).

For the last assertion we have to show that there are no other compact objects. Since constructibility and finite Tor amplitude are preserved by direct summands, the category \(D^b_{\mathrm{ctf}}(X,A)\) is thick. It is essentially small, because it is a full subcategory of the essentially small bounded constructible category, and by definition a skeleton generates \(\Loc(D^b_{\mathrm{ctf}}(X,A))\). In a triangulated category with direct sums generated by a set of compact objects, the compact objects form the thick closure of that set [Stacks, \href{https://stacks.math.columbia.edu/tag/09SM}{09SM}, \href{https://stacks.math.columbia.edu/tag/09SN}{09SN}]. In our case this closure is \(D^b_{\mathrm{ctf}}(X,A)\).
\end{proof}

% text: prose-05-17
\phantomsection\label{bp47}%
For finite coefficients the localizing category is a familiar one: the next proposition shows that it is the classical unbounded derived category of the étale site. In particular, \(f^!\) with finite coefficients is defined on all of \(D(X_{\acute{e}t},A)\). The proposition also shows that \(Rf_!\), once it is defined as \(Rp_*j_!\) by a compactification in Section \hyperref[bp53]{5.4}, is computed by the classical functor \(Rp_{\acute{e}t,*}j_{\acute{e}t,!}\) on the étale site.

% object: proposition-comparison-etale-derived-category (blueprint/chap-Construction-of-f/sec-Compactness-for-finite-coefficients/proposition-comparison-etale-derived-category/_section.md)
\phantomsection\label{bp48}%
\begin{bpstatementit}{Proposition~5.15}{ (Comparison with the étale derived category)}
For a finite admissible ring \(A\), the comparison functor \(\nu_X^*:D(X_{\acute{e}t},A)\to D(X_{\mathrm{pro}\acute{e}t},A)\) identifies

\[
D(X_{\acute{e}t},A)=\Loc(D^b_{\mathrm{ctf}}(X_{\acute{e}t},A))
\simeq\Loc(D^b_{\mathrm{ctf}}(X,A)).
\]

It preserves arbitrary direct sums. For every morphism \(g:X\to Y\) of schemes of finite type over \(k\) and every \(K\in D(X_{\acute{e}t},A)\), the comparison map (2.9) \(\nu_Y^*Rg_{\acute{e}t,*}K\to Rg_{\mathrm{pro}\acute{e}t,*}\nu_X^*K\) is an isomorphism, and \(\nu^*\) commutes with inverse images and, by (2.10), with extension by zero along open immersions.
\end{bpstatementit}

% proof: proof/proposition-comparison-etale-derived-category
\begin{proof}[{Proof}]\phantomsection\label{bp49}%
There are two statements: \(\nu_X^*\) identifies the two localizing categories, and on the étale side the localizing category is everything. The first follows formally from the comparison results for the pro-étale site. For the second we exhibit compact generators of \(D(X_{\acute{e}t},A)\) that lie in \(D^b_{\mathrm{ctf}}\).

An exact fully faithful functor that preserves direct sums identifies the localizing subcategories generated by corresponding generators. In our case, Proposition \bpout{2.49} shows that \(\nu_X^*\) is such a functor on unbounded complexes and that it is compatible with direct images, and on \(D^b_{\mathrm{ctf}}\) it is the equivalence of Corollary \bpout{2.44}. Finally, inverse images commute with \(\nu^*\) because both composites are inverse images along the same morphism of topoi.

For the equality with the whole étale derived category, let \(u:U\to X\) range over a basis of quasi-compact separated étale objects. The free representable module sheaves \(u_!A_U\) are flat constructible, hence belong to \(D^b_{\mathrm{ctf}}(X_{\acute{e}t},A)\). We have \(R\!\operatorname{Hom}_X(u_!A_U,K)=R\Gamma(U,K)\). It follows that their shifts detect zero objects of \(D(X_{\acute{e}t},A)\). Indeed, if \(H^n(U,K)=0\) for all such \(U\) and all \(n\), then every cohomology sheaf of \(K\) vanishes, being the sheafification of \(U\mapsto H^n(U,K)\). These objects are also compact in all of \(D(X_{\acute{e}t},A)\). The reason is that \(R\Gamma(U,-)\) commutes with direct sums, since \(\operatorname{cd}_\ell(U)<\infty\), as in the étale part of the proof of Proposition \hyperref[bp45]{5.14}. A set of compact objects whose shifts detect zero objects generates a triangulated category with direct sums as a localizing subcategory. Hence these objects generate the entire étale derived category.
\end{proof}

% text: prose-05-18
\phantomsection\label{bp50}%
Proper direct image already preserves \(D^b_{\mathrm{ctf}}\) by Proposition \bpout{4.35}. Its extension to the localizing categories therefore only needs preservation of arbitrary direct sums. By the comparison just proved, this reduces to the unbounded étale statement.

% object: corollary-proper-direct-image-finite-coefficients-localizing (blueprint/chap-Construction-of-f/sec-Compactness-for-finite-coefficients/corollary-proper-direct-image-finite-coefficients-localizing/_section.md)
\phantomsection\label{bp51}%
\begin{bpstatementit}{Corollary~5.16}{ (Proper direct image on the localizing categories)}
Let \(A\) be a finite admissible ring and let \(p:X\to Y\) be proper. The functor \(Rp_*\) restricts to an exact functor

\[
Rp_*:\Loc(D^b_{\mathrm{ctf}}(X,A))\longrightarrow\Loc(D^b_{\mathrm{ctf}}(Y,A))
\]

that preserves arbitrary direct sums.
\end{bpstatementit}

% proof: proof/corollary-proper-direct-image-finite-coefficients-localizing
\begin{proof}[{Proof}]\phantomsection\label{bp52}%
On unbounded étale derived categories, proper direct image preserves arbitrary direct sums by [Stacks, \href{https://stacks.math.columbia.edu/tag/0F0D}{0F0D}]. Forgetting coefficients gives the \(A\)-linear statement, and Proposition \hyperref[bp48]{5.15} transports it to the pro-étale localizing categories.

Since \(Rp_*\) maps \(D^b_{\mathrm{ctf}}(X,A)\) into \(D^b_{\mathrm{ctf}}(Y,A)\) by Proposition \bpout{4.35}, the objects \(K\in\Loc(D^b_{\mathrm{ctf}}(X,A))\) whose direct image belongs to \(\Loc(D^b_{\mathrm{ctf}}(Y,A))\) form a triangulated subcategory closed under direct sums and containing \(D^b_{\mathrm{ctf}}(X,A)\), that is, all of \(\Loc(D^b_{\mathrm{ctf}}(X,A))\).
\end{proof}


% object: sec-Direct-image-with-compact-support (blueprint/chap-Construction-of-f/sec-Direct-image-with-compact-support/_section.md), content hidden
\subsection*{\texorpdfstring{5.4\quad Direct image with compact support}{5.4 Direct image with compact support}}
\phantomsection\label{bp53}%
\addcontentsline{toc}{subsection}{\texorpdfstring{5.4\quad Direct image with compact support}{5.4 Direct image with compact support}}
% text: prose-05-19
\phantomsection\label{bp54}%
The trace formula needs a relative version of cohomology with compact supports, whose stalks are the groups \(H^i_c\) of the fibres. Consider first a locally compact Hausdorff space \(X\) with an open embedding \(j:X\hookrightarrow\overline X\) into a compact space. A section of \(j_!F\) is a section of \(F\) that vanishes near \(\overline X\setminus X\), and \(H^i_c(X,F)=H^i(\overline X,j_!F)\). For example, for \(X=\mathbf R\) and \(\overline X\) the circle, we find \(H^1_c(\mathbf R,\mathbf Z)=\mathbf Z\).

The definition by sections with compact support has no direct analogue for schemes, because every scheme of finite type over \(k\) is quasi-compact, and the algebraic substitute for compactness is properness, which is a property of a morphism and not of a subset. The two-step form, however, does carry over. In it, the functor \(j_!\) imposes compact support along the open immersion, while for a proper map \(p\), as for a map of spaces with compact fibres, the direct image with compact supports is the ordinary direct image \(Rp_*\). By proper base change, the result does not depend on the compactification (Proposition \hyperref[bp86]{5.26}). Proper base change also identifies the stalks of \(Rf_!\) with the compactly supported cohomology of the fibres (Corollary \bpout{12.26}).

Let \(f:X\to Y\) be a separated morphism of finite type. For a compactification \(f=p\circ j\), with \(j\) an open immersion and \(p\) proper, we define \(Rf_!=Rp_*\circ j_!\). We must show that this functor preserves direct sums. In the following subsection we take \(\Lambda=E\) or \(\overline{\mathbf Q}_\ell\), so that \(\Loc(D^b_{\mathrm{ctf}}(X,\Lambda))=\Loc(\Dbc(X,\Lambda))\).

% object: subsec-Direct-sums (blueprint/chap-Construction-of-f/sec-Direct-image-with-compact-support/subsec-Direct-sums/_section.md), content hidden
\subsubsection*{\texorpdfstring{5.4.1\quad Direct sums}{5.4.1 Direct sums}}
\phantomsection\label{bp55}%
\addcontentsline{toc}{subsubsection}{\texorpdfstring{5.4.1\quad Direct sums}{5.4.1 Direct sums}}
% text: prose-05-20
\phantomsection\label{bp56}%
To apply Brown representability, we need \(Rf_!\) to preserve arbitrary direct sums. Since \(j_!\) is a left adjoint, it does. The question is therefore whether \(Rp_*\) preserves direct sums for proper \(p\). For complexes with a common lower and upper bound, this would follow from finiteness of \(\operatorname{cd}_\ell(p)\). However, the summands of a direct sum in \(\Loc(D^b_{\mathrm{ctf}}(X,\Lambda))\) need not have any common bound. As in the proof of compactness, we use a uniform bound instead. Let \(d\) bound the fibre dimensions of \(p\). We show that \(R^qp_*\) vanishes on \(\IndCons(X,\Lambda)\) for \(q>2d+1\). This means that \(\mathcal H^q(Rp_*K)\) depends only on the cohomology of \(K\) in degrees \([q-2d-1,q]\).

We prove the bound on sections over affine weakly contractible objects and then pass to colimits. For this we use the continuity of pro-étale cohomology, Corollary \bpout{2.29}, applied to sheaves, which are complexes vanishing below degree \(0\). It says that for a quasi-compact and quasi-separated morphism \(p\) the functors \(R^qp_*\) on sheaves of modules commute with filtered colimits. Since a direct sum is the filtered colimit of its finite subsums, they also commute with direct sums.

We also use the following description of \(R^qp_*\). If \(V\to Y\) is an affine weakly contractible object and \(X_V=X\times_YV\), then \(\Gamma(V,-)\) is exact (Corollary corollary-evaluation-weakly-contractible-objects), and hence

\[
(R^qp_*M)(V)=H^q((X_V)_{\mathrm{pro}\acute{e}t},M|_{X_V}).
\tag{5.5}
\]

Since such \(V\) form a basis of \(Y_{\mathrm{pro}\acute{e}t}\), a sheaf on \(Y_{\mathrm{pro}\acute{e}t}\) vanishes if its sections over all of them vanish.

% object: lemma-uniform-bound-proper-direct-image (blueprint/chap-Construction-of-f/sec-Direct-image-with-compact-support/subsec-Direct-sums/lemma-uniform-bound-proper-direct-image/_section.md)
\phantomsection\label{bp57}%
\begin{bpstatementit}{Lemma~5.17}{ (A uniform bound for proper direct image)}
Let \(p:X\to Y\) be proper, and let \(d\) be an integer bounding the dimensions of all geometric fibres of \(p\) (for instance \(d=\dim X\)). For \(M\in\IndCons(X,\Lambda)\),

\[
R^qp_*M=0\quad(q>2d+1),\qquad R^qp_*M\in\IndCons(Y,\Lambda)\quad(q\geq0),
\]

and \(R^qp_*:\IndCons(X,\Lambda)\to\IndCons(Y,\Lambda)\) preserves arbitrary direct sums.
\end{bpstatementit}

% text: prose-05-21
\phantomsection\label{bp58}%
By a stratification we reduce the lemma to the extension by zero of a lisse \(\mathcal O_E\)-lattice on a stratum. Its reductions modulo \(\pi_E^n\) are torsion sheaves, for which the classical bound \(2d\) holds. In the passage to the limit over \(n\), the \(\varprojlim^1\) term raises the bound by one, to \(2d+1\). The remaining steps, inverting \(\ell\), extending scalars to \(\overline{\mathbf Q}_\ell\) and passing to ind-objects, are filtered colimits. By Corollary \bpout{2.29}, the bound survives each of them.

% proof: proof/lemma-uniform-bound-proper-direct-image
\begin{proof}[{Proof of Lemma 5.17}]\phantomsection\label{bp59}%
\emph{Reduction to a lattice.} The vanishing statement is stable under extensions: if \(R^qp_*\) vanishes for \(q>2d+1\) on two sheaves, it vanishes on every extension of one by the other. By Lemma lemma-constructible-sheaves, every \(M\in\Cons(X,\Lambda)\) is a finite extension of sheaves \(a_!L\), with \(a:S\hookrightarrow X\) a locally closed immersion and \(L\) lisse. Moreover, for \(\Lambda=\overline{\mathbf Q}_\ell\), \(L\) is the scalar extension \(\QbarX\otimes_{E}L_E\) of a lisse \(E_S\)-module \(L_E\) for some finite \(E/\mathbf Q_\ell\) (Lemma \bpout{3.17}). For \(\Lambda=E\), we simply put \(L_E=L\). In both cases we may refine \(S\) by a finite stratification, since this replaces \(a_!L\) by a finite extension of sheaves of the same form. By Lemma \bpout{3.24}, after such a refinement \(L_E\) contains a lattice \(L_0\), locally free of finite rank over \(\mathcal O_{E,S}\), with \(L_0[1/\ell]=L_E\). Hence it is enough to prove the bound for \(a_!L_0\), and then to pass to \(a_!L_E=(a_!L_0)[1/\ell]\) and, if \(\Lambda=\overline{\mathbf Q}_\ell\), to its scalar extension \(a_!L\).

In the passage to the limit below we use that the sheaf \(a_!L_0\) is derived \(\pi_E\)-complete. To see this, note that \(a\) is a closed immersion composed with a quasi-compact open immersion. Closed pushforward is exact and preserves products [Stacks, \href{https://stacks.math.columbia.edu/tag/09BM}{09BM}], and extension by zero is exact and preserves limits [Stacks, \href{https://stacks.math.columbia.edu/tag/09BP}{09BP}]. Therefore both functors preserve the derived inverse limit that expresses the completeness of \(L_0\).

\emph{The torsion bound.} Over the objects \(V\) that compute \(R^qp_*\), the reductions \(a_!(L_0/\pi_E^n)\) have cohomology only up to degree \(2d\). Fix an affine weakly contractible \(V\to Y\) and let \(p_V:X_V\to V\) be the base change of \(p\), which is proper with the same bound \(d\). Since \(L_0\) is flat, \(a_!L_0\otimes^{\mathbf L}_{\mathcal O_E}
\mathcal O_E/\pi_E^n\) is the sheaf \(a_!(L_0/\pi_E^n)\), which is the pullback to the pro-étale site of a classical constructible étale sheaf of \(\mathcal O_E/\pi_E^n\)-modules. Since the scheme \(V\) is not Noetherian, we use the proper vanishing theorem in the generality of [Stacks, \href{https://stacks.math.columbia.edu/tag/095U}{095U}]: if a proper morphism has geometric fibres of dimension at most \(d\), then its higher direct images of torsion sheaves vanish in degrees \(q>2d\). With the comparison [Stacks, \href{https://stacks.math.columbia.edu/tag/09A6}{09A6}] of étale and pro-étale direct images, this theorem for \(p_V\) says that \(R^qp_{V,*}\) of this sheaf vanishes for \(q>2d\). As \(V\) is weakly contractible, \(\Gamma(V,-)\) is exact on all pro-étale sheaves [Stacks, \href{https://stacks.math.columbia.edu/tag/0991}{0991}]. Hence \(H^q(X_V,-)=\Gamma(V,R^qp_{V,*}(-))\), and

\[
H^q\bigl(X_V,a_!(L_0/\pi_E^n)\bigr)=0\quad(q>2d).
\]

\emph{The limit.} Derived completeness lets us pass from the reductions to \(a_!L_0\). Since restriction to \(X_V\) is exact and commutes with products, hence with derived inverse limits, completeness gives \((a_!L_0)|_{X_V}\simeq R\varprojlim_n a_!(L_0/\pi_E^n)|_{X_V}\). Derived global sections, too, commute with derived inverse limits, since a product of K-injective complexes is K-injective and sections preserve products. The Milnor exact sequence of Proposition \bpout{2.24}

\[
0\to{\textstyle\varprojlim^1_n}H^{q-1}\bigl(X_V,a_!(L_0/\pi_E^n)\bigr)
\to H^q(X_V,a_!L_0)
\to\varprojlim_nH^q\bigl(X_V,a_!(L_0/\pi_E^n)\bigr)\to0
\]

then shows that \(H^q(X_V,a_!L_0)=0\) for \(q>2d+1\), since for such \(q\) both outer terms vanish by the torsion bound.

\emph{Back to \(\Lambda\).} The remaining passages are filtered colimits. By Corollary \bpout{2.29}, cohomology on the quasi-compact and quasi-separated scheme \(X_V\) commutes with them. For instance, the sheaf \((a_!L_0)[1/\ell]\) is the filtered colimit of \(a_!L_0\) along multiplication by \(\ell\), and so it has the same vanishing. For \(\Lambda=\overline{\mathbf Q}_\ell\), scalar extension from \(E\) to \(\overline{\mathbf Q}_\ell\) is the filtered colimit of the extensions to finite \(E'\supset E\). Since \(E'\) is finite free over \(E\), each of these is, as a sheaf of abelian groups, a finite direct sum of copies of the original sheaf. Hence the continuity lemma again gives the bound. As \(V\) was an arbitrary affine weakly contractible object, formula (5.5) shows that \(R^qp_*(a_!L)=0\) for \(q>2d+1\). By the first step, \(R^qp_*M=0\) for \(q>2d+1\) and every \(M\in\Cons(X,\Lambda)\).

\emph{Passage to \(\IndCons(X,\Lambda)\).} For constructible \(M\) we have \(R^qp_*M\in\Cons(Y,\Lambda)\), because \(Rp_*\) maps \(\Dbc(X,\Lambda)\) into \(\Dbc(Y,\Lambda)\) by Proposition \bpout{4.40}. Moreover, every object of \(\IndCons(X,\Lambda)\) is a filtered colimit of constructible sheaves, and \(R^qp_*\) commutes with filtered colimits and direct sums by Corollary \bpout{2.29}. All three assertions for \(\IndCons(X,\Lambda)\) therefore follow.
\end{proof}

% text: prose-05-22
\phantomsection\label{bp60}%
The bound \(2d+1\) is not sharp. For example, for the structure map of a smooth projective curve over \(\bar k\), where \(d=1\), the lemma only says that cohomology vanishes above degree \(3\), while in fact it vanishes above degree \(2\). The extra degree does no harm, because the arguments below only need a bound that does not depend on the sheaf. As in the compactness argument, such a bound passes from sheaves to unbounded complexes.

% object: lemma-proper-direct-image-unbounded-complexes (blueprint/chap-Construction-of-f/sec-Direct-image-with-compact-support/subsec-Direct-sums/lemma-proper-direct-image-unbounded-complexes/_section.md)
\phantomsection\label{bp61}%
\begin{bpstatementit}{Lemma~5.18}{ (Proper direct image of unbounded complexes)}
Let \(p:X\to Y\) be proper, and let \(d\) bound the dimensions of the geometric fibres of \(p\). For integers \(a\) and \(b\), write \(\Dic^{\leq b}(X,\Lambda)\) and \(\Dic^{\geq a}(X,\Lambda)\) for the objects of \(\Dic(X,\Lambda)\) (Corollary \hyperref[bp33]{5.11}) with cohomology in degrees \(\leq b\), respectively \(\geq a\).

\begin{enumerate}
\item \(Rp_*\) maps \(\Dic(X,\Lambda)\) into \(\Dic(Y,\Lambda)\), with \(Rp_*\Dic^{\leq b}(X,\Lambda)\subset\Dic^{\leq b+2d+1}(Y,\Lambda)\) and \(Rp_*\Dic^{\geq a}(X,\Lambda)\subset\Dic^{\geq a}(Y,\Lambda)\).
\item For \(K\in\Dic(X,\Lambda)\) and every integer \(q\), the truncation maps \(K\leftarrow\tau_{\leq q}K\to\tau_{\geq q-2d-1}\tau_{\leq q}K\) induce an isomorphism in \(\IndCons(Y,\Lambda)\) \[
   \mathcal H^q(Rp_*K)\simeq
   \mathcal H^q\bigl(Rp_*\tau_{\geq q-2d-1}\tau_{\leq q}K\bigr).
   \tag{5.6}
   \]
\item \(Rp_*:\Dic(X,\Lambda)\to\Dic(Y,\Lambda)\) preserves arbitrary direct sums.
\end{enumerate}
\end{bpstatementit}

% proof: proof/lemma-proper-direct-image-unbounded-complexes
\begin{proof}[{Proof}]\phantomsection\label{bp62}%
For bounded below \(K\) everything is read off from the hypercohomology spectral sequence, which has only \(2d+2\) nonzero columns. A general \(K\), in turn, is the derived limit of its truncations \(\tau_{\geq-n}K\), and the column bound shows that in each fixed degree the tower of direct images is eventually constant. The window (5.6) and the preservation of direct sums follow from this.

\emph{Bounded below \(K\).} The hypercohomology spectral sequence

\[
E_2^{u,v}=R^up_*\mathcal H^v(K)\Longrightarrow\mathcal H^{u+v}(Rp_*K)
\]

converges. By Lemma \hyperref[bp57]{5.17}, its only nonzero columns are \(0\leq u\leq2d+1\), and their terms lie in \(\IndCons(Y,\Lambda)\). So each \(\mathcal H^n(Rp_*K)\) has a finite filtration whose graded pieces are subquotients of objects of \(\IndCons(Y,\Lambda)\). Since this category is closed under kernels, cokernels and extensions by Proposition \hyperref[bp27]{5.9}, we get \(\mathcal H^n(Rp_*K)\in\IndCons(Y,\Lambda)\). Both amplitude bounds can be read off from the columns.

\emph{Arbitrary \(K\).} By Proposition \bpout{2.26}, \(K\simeq R\varprojlim_n
\tau_{\geq-n}K\). Since \(Rp_*\) is a right adjoint, it preserves products and hence derived inverse limits. Thus \(Rp_*K\simeq R\varprojlim_nRp_*\tau_{\geq-n}K\), and it suffices to see that in each fixed degree the tower \(Rp_*\tau_{\geq-n}K\) stabilizes. The fibre of \(\tau_{\geq-n-1}K\to\tau_{\geq-n}K\) is \(\mathcal H^{-n-1}(K)[n+1]\), a single ind-constructible sheaf in degree \(-n-1\). By the bounded below case its direct image has cohomology only in degrees \([-n-1,-n-1+2d+1]=[-n-1,2d-n]\). In particular it has none in degrees \(q-1\), \(q\), \(q+1\) once \(n>2d+1-q\). For fixed \(q\), the long exact sequences then show that the transition maps of the towers \(\mathcal H^q(Rp_*\tau_{\geq-n}K)\) and \(\mathcal H^{q-1}(Rp_*\tau_{\geq-n}K)\) are isomorphisms for \(n\) large. Since an eventually constant tower has vanishing \(\varprojlim^1\), the Milnor sequence of Proposition \bpout{2.24} gives

\[
\mathcal H^q(Rp_*K)\simeq
\mathcal H^q(Rp_*\tau_{\geq-n}K)\quad\text{for all sufficiently large }n.
\]

The bounded below case, applied to the right-hand side, shows that \(\mathcal H^q(Rp_*K)\in\IndCons(Y,\Lambda)\), with the bound \(b+2d+1\) if \(K\in
\Dic^{\leq b}(X,\Lambda)\). The lower bound needs no new argument, because an object of \(\Dic^{\geq a}\) is bounded below.

\emph{The isomorphism (5.6).} It comes from the bounds, applied to the truncation triangles

\[
\tau_{\leq q}K\to K\to\tau_{\geq q+1}K,
\qquad
\tau_{\leq q-2d-2}K\to\tau_{\leq q}K\to\tau_{\geq q-2d-1}\tau_{\leq q}K.
\]

The direct image of \(\tau_{\geq q+1}K\) lives in degrees \(\geq q+1\), and that of \(\tau_{\leq q-2d-2}K\) in degrees \(\leq q-1\). The long exact cohomology sequences of the two triangles therefore yield isomorphisms \(\mathcal H^q(Rp_*\tau_{\leq q}K)\simeq\mathcal H^q(Rp_*K)\) and \(\mathcal H^q(Rp_*\tau_{\leq q}K)\simeq
\mathcal H^q(Rp_*\tau_{\geq q-2d-1}\tau_{\leq q}K)\).

\emph{Direct sums, assertion (3).} The window (5.6) reduces us to families with cohomology in a common finite interval, and for such a family \((M_\alpha)\) we show by induction on the length of the interval that the map \(\bigoplus_\alpha Rp_*M_\alpha\to Rp_*(\bigoplus_\alpha M_\alpha)\) is an isomorphism. When the length is zero, each \(M_\alpha\) is a shifted object of \(\IndCons(X,\Lambda)\), and the claim is the assertion about direct sums in Lemma \hyperref[bp57]{5.17}, in each degree. Then, for the inductive step, we take the direct sum of the truncation triangles that split off the top cohomology sheaf and apply the five lemma to the long exact cohomology sequences.

Now let \((K_\alpha)\) be arbitrary and fix \(q\). Since direct sums of sheaves are exact, cohomology and truncations commute with direct sums. The objects \(\tau_{\geq q-2d-1}\tau_{\leq q}K_\alpha\) have cohomology in the common interval \([q-2d-1,q]\). Applying (5.6) to each \(K_\alpha\) and to \(\bigoplus_\alpha K_\alpha\), and the first paragraph to the truncated family, we see that \(\bigoplus_\alpha Rp_*K_\alpha\to Rp_*(\bigoplus_\alpha K_\alpha)\) is an isomorphism on \(\mathcal H^q\). Since \(q\) is arbitrary, the map is an isomorphism.
\end{proof}


% object: subsec-Integral-coefficients (blueprint/chap-Construction-of-f/sec-Direct-image-with-compact-support/subsec-Integral-coefficients/_section.md), content hidden
\subsubsection*{\texorpdfstring{5.4.2\quad Integral coefficients}{5.4.2 Integral coefficients}}
\phantomsection\label{bp63}%
\addcontentsline{toc}{subsubsection}{\texorpdfstring{5.4.2\quad Integral coefficients}{5.4.2 Integral coefficients}}
% text: prose-05-23
\phantomsection\label{bp64}%
For \(\mathcal O_E\) the rational arguments for compactness (Theorem \hyperref[bp40]{5.13}) and for direct sums (Lemma \hyperref[bp61]{5.18}) go through with one change. Integral constructible sheaves are no longer locally free on strata; for example, \(\mathcal O_{E,X}/\pi_E\) is constructible. But \(\mathcal O_E\) is a discrete valuation ring, so they have free resolutions of length one. The estimates only use the resulting Tor amplitude \([-1,0]\).

% object: lemma-structure-integral-constructible-complexes (blueprint/chap-Construction-of-f/sec-Direct-image-with-compact-support/subsec-Integral-coefficients/lemma-structure-integral-constructible-complexes/_section.md)
\phantomsection\label{bp65}%
\begin{bpstatementit}{Lemma~5.19}{ (Structure of integral constructible complexes)}
\begin{enumerate}
\item \(\Dbc(X,\mathcal O_E)\) is a triangulated subcategory of \(D(X_{\mathrm{pro}\acute{e}t},\mathcal O_E)\) of bounded objects, stable under truncation; its objects in degree \(0\) form the category \(\Cons(X,\mathcal O_E)\), an essentially small abelian subcategory of \(\operatorname{Mod}(X_{\mathrm{pro}\acute{e}t},\mathcal O_E)\) closed under extensions.
\item For \(C\in\Cons(X,\mathcal O_E)\) there is a finite stratification on whose strata \(C\) is pro-étale locally \(\mathcal O_{E,S}\otimes_{\mathcal O_E}P\), with \(P\) finitely generated.
\item On each such stratum \(C\) is dualizable of Tor amplitude \([-1,0]\), and its dual has Tor amplitude \([0,1]\).
\end{enumerate}
\end{bpstatementit}

% proof: proof/lemma-structure-integral-constructible-complexes
\begin{proof}[{Proof}]\phantomsection\label{bp66}%
The conditions of Proposition \bpout{3.11}, derived completeness and constructible reduction, are stable under shifts and cones, so \(\Dbc(X,\mathcal O_E)\) is triangulated. By Proposition \bpout{3.10}, on the strata of a finite stratification an object is pro-étale locally of the form \(\mathcal O_{E,S}\otimes P^\bullet\) with \(P^\bullet\) a bounded complex of finite free \(\mathcal O_E\)-modules. Since \(\mathcal O_{E,S}\) is flat over \(\mathcal O_E\), we get \(H^i=\mathcal O_{E,S}\otimes H^i(P^\bullet)\) on such a stratum. This local form of the cohomology sheaves shows that the truncations stay in \(\Dbc\). It also gives (2), and then the abelian part of (1). In (1), essential smallness holds because a constructible sheaf is determined up to isomorphism by a stratification, local presentations, descent data and finitely many extension classes. Finally, for (3) we use that a finitely generated module over the discrete valuation ring \(\mathcal O_E\) has a free resolution \(0\to\mathcal O_E^r\to\mathcal O_E^s\to P\to0\) of length one.
\end{proof}

% text: prose-05-24
\phantomsection\label{bp67}%
For instance, \(\mathcal O_{E,X}/\pi_E\) is the cokernel of \(\pi_E:\mathcal O_{E,X}\to\mathcal O_{E,X}\), so it has Tor amplitude \([-1,0]\). Its dual \(\cHom(\mathcal O_{E,X}/\pi_E,\mathcal O_{E,X})=
(\mathcal O_{E,X}/\pi_E)[-1]\) has Tor amplitude \([0,1]\), as in (3).

The dévissage in the proof of compactness also needs that the operations it uses preserve \(\Dbc(-,\mathcal O_E)\).

% object: lemma-inverse-image-tensor-product (blueprint/chap-Construction-of-f/sec-Direct-image-with-compact-support/subsec-Integral-coefficients/lemma-inverse-image-tensor-product/_section.md)
\phantomsection\label{bp68}%
\begin{bpstatementit}{Lemma~5.20}{ (Inverse image and tensor product)}
Inverse image, \(\otimes^{\mathbf L}_{\mathcal O_{E,X}}\) without completion, extension by zero along a quasi-compact open immersion and direct image along a closed immersion preserve \(\Dbc(-,\mathcal O_E)\); on \(\Dbc\) inverse image agrees with the completed inverse image of the Stacks Project.
\end{bpstatementit}

% proof: proof/lemma-inverse-image-tensor-product
\begin{proof}[{Proof}]\phantomsection\label{bp69}%
The preservation statements follow from Lemma lemma-integral-models, since closed direct image is exact and preserves limits [Stacks, \href{https://stacks.math.columbia.edu/tag/09AA}{09AA}, \href{https://stacks.math.columbia.edu/tag/09BM}{09BM}]. The two inverse images agree on objects that are pro-étale locally of the form \(\mathcal O_{E,S}\otimes P\) with \(P\) finitely generated, and by the local structure of Lemma lemma-structure-integral-constructible-complexes and (2) they therefore agree on \(\Dbc\).
\end{proof}

% text: prose-05-25
\phantomsection\label{bp70}%
With these two lemmas, the integral case of the two hypotheses of Brown representability can be checked by following the rational proofs.

% object: proposition-compactness-proper-direct-image (blueprint/chap-Construction-of-f/sec-Direct-image-with-compact-support/subsec-Integral-coefficients/proposition-compactness-proper-direct-image/_section.md)
\phantomsection\label{bp71}%
\begin{bpstatementit}{Proposition~5.21}{ (Compactness and proper direct image)}
Every object of \(D^b_{\mathrm{ctf}}(X,\mathcal O_E)=\Dbc(X,\mathcal O_E)\) is compact in \(\Loc(D^b_{\mathrm{ctf}}(X,\mathcal O_E))\), and for proper \(p:X\to Y\) the functor \(Rp_*\) preserves direct sums of objects of \(\Loc(D^b_{\mathrm{ctf}}(X,\mathcal O_E))\).
\end{bpstatementit}

% proof: proof/proposition-compactness-proper-direct-image
\begin{proof}[{Proof}]\phantomsection\label{bp72}%
The rational proofs (Theorem \hyperref[bp40]{5.13}, Lemmas \hyperref[bp37]{5.12} and \hyperref[bp61]{Lemma 5.18}) go through once we supply for \(\mathcal O_E\) the four inputs they use: the abelian structure of Lemma lemma-structure-integral-constructible-complexes, continuity of \(\Ext^i(C,-)\) for bounded below systems, a bound \(d_C\) with \(\Ext^i(C,N)=0\) for \(i>d_C\) and \(N\) constructible, and a bound on the amplitude of \(Rp_*\).

For continuity, take \(C=a_!L\) on a stratum. By Lemma lemma-sections-support-closed-subscheme, \(R\!\operatorname{Hom}(a_!L,-)\) is \(R\Gamma(S,-)\) of the dual of \(L\) tensored with \(u^*R\underline\Gamma_Z\). The proof of that formula uses only that \(L\) is dualizable, which holds by Lemma lemma-structure-integral-constructible-complexes, and so continuity follows from Corollary \bpout{2.29}.

The bound \(d_C\) comes from counting degrees as in Lemma \hyperref[bp20]{5.6}. Namely, the reductions of \(N\) lie in \([-1,0]\), and the dual of \(L\) adds one degree. Next, the functors \(R\Gamma\) and \(R\underline\Gamma_Z\) add the cohomological dimensions of Definition \bpout{2.45}. Finally, the derived limit over the reductions adds one more [Stacks, \href{https://stacks.math.columbia.edu/tag/08TB}{08TB}].

The same count bounds \(Rp_*\). For fibres of dimension \(\le d\), the torsion bound [Stacks, \href{https://stacks.math.columbia.edu/tag/095U}{095U}] gives \(2d\), and the derived limit adds one. Hence \(Rp_*\) of a constructible sheaf lies in \([0,2d+1]\).
\end{proof}

% text: prose-05-26
\phantomsection\label{bp73}%
These are the facts about \(\mathcal O_E\) that we use in the construction of \(Rf_!\) below.


% object: subsec-Construction-of-Rf (blueprint/chap-Construction-of-f/sec-Direct-image-with-compact-support/subsec-Construction-of-Rf/_section.md), content hidden
\subsubsection*{\texorpdfstring{5.4.3\quad Construction of $Rf_!$}{5.4.3 Construction of Rf!}}
\phantomsection\label{bp74}%
\addcontentsline{toc}{subsubsection}{\texorpdfstring{5.4.3\quad Construction of $Rf_!$}{5.4.3 Construction of Rf!}}
% text: prose-05-27
\phantomsection\label{bp75}%
From now on, \(\Lambda\) is any admissible coefficient ring. Lemmas \hyperref[bp57]{5.17} and \hyperref[bp61]{Lemma 5.18} only treat \(E\) and \(\overline{\mathbf Q}_\ell\), but we use them only through the fact that \(Rp_*\) preserves direct sums for proper \(p\). For finite and integral coefficients this fact was proved separately, and the next proof collects the references.

Proper base change shows, by one and the same argument, that \(Rp_*\circ j_!\) does not depend on the compactification and is compatible with composition and base change. The comparison maps are maps out of an extension by zero \(j_!K\). They are therefore determined by their restriction to \(X\), where they are the identity. On the closed complement of \(X\), proper base change shows that both sides vanish, and so the map is an isomorphism. For instance, if \(h:Z'\to Z\) is proper and \(h^{-1}(X)=X\), then \(Rh_*j'_!K\) restricted to \(Z\setminus X\) is the direct image of the restriction of \(j'_!K\) to \(Z'\setminus X\), which is zero. Since proper base change is only known on \(D^b_{\mathrm{ctf}}\), we first show that all the functors involved preserve direct sums, so that it suffices to check isomorphisms on \(D^b_{\mathrm{ctf}}\) (Lemma \hyperref[bp79]{5.23}).

That both factors of \(Rp_*\circ j_!\) preserve the localizing categories comes from the definition of \(\Loc(D^b_{\mathrm{ctf}}(X,\Lambda))\) in the following form: if the objects of \(\Loc(D^b_{\mathrm{ctf}}(X,\Lambda))\) with some property form a triangulated subcategory that is closed under direct sums and contains \(D^b_{\mathrm{ctf}}(X,\Lambda)\), then every object of \(\Loc(D^b_{\mathrm{ctf}}(X,\Lambda))\) has the property.

% object: proposition-proper-direct-image-extension-zero (blueprint/chap-Construction-of-f/sec-Direct-image-with-compact-support/subsec-Construction-of-Rf/proposition-proper-direct-image-extension-zero/_section.md)
\phantomsection\label{bp76}%
\begin{bpstatementit}{Proposition~5.22}{ (Proper direct image and extension by zero)}
For every admissible \(\Lambda\), if \(p:X\to Y\) is proper, \(Rp_*:D(X_{\mathrm{pro}\acute{e}t},\Lambda)\to D(Y_{\mathrm{pro}\acute{e}t},\Lambda)\) restricts to an exact functor \(\Loc(D^b_{\mathrm{ctf}}(X,\Lambda))\to\Loc(D^b_{\mathrm{ctf}}(Y,\Lambda))\) that preserves direct sums. If \(j:U\to X\) is an open immersion, \(j_!:D(U_{\mathrm{pro}\acute{e}t},\Lambda)\to D(X_{\mathrm{pro}\acute{e}t},\Lambda)\) restricts to an exact functor \(\Loc(D^b_{\mathrm{ctf}}(U,\Lambda))\to\Loc(D^b_{\mathrm{ctf}}(X,\Lambda))\) that preserves direct sums and maps \(D^b_{\mathrm{ctf}}(U,\Lambda)\) into \(D^b_{\mathrm{ctf}}(X,\Lambda)\).
\end{bpstatementit}

% proof: proof/proposition-proper-direct-image-extension-zero
\begin{proof}[{Proof}]\phantomsection\label{bp77}%
For each functor we need two facts: it preserves direct sums, and it maps \(D^b_{\mathrm{ctf}}\) into \(D^b_{\mathrm{ctf}}\). The rest then follows from the definition of the localizing categories, in the form recalled before the proposition.

For \(Rp_*\) with \(E\) and \(\overline{\mathbf Q}_\ell\), Lemma lemma-proper-direct-image-unbounded-complexes shows that \(Rp_*\) preserves direct sums on \(\Dic(X,\Lambda)\supset\Loc(D^b_{\mathrm{ctf}}(X,\Lambda))\), and Proposition \bpout{4.40} shows that \(Rp_*\) preserves \(D^b_{\mathrm{ctf}}=\Dbc\). For finite \(A\), these two facts are Proposition \bpout{4.35} and Corollary \hyperref[bp51]{5.16}; for \(\mathcal O_E\), they are Proposition \hyperref[bp71]{5.21} and Proposition \bpout{4.40}. In each case, the objects \(K\in\Loc(D^b_{\mathrm{ctf}}(X,\Lambda))\) with \(Rp_*K\in\Loc(D^b_{\mathrm{ctf}}(Y,\Lambda))\) form a triangulated subcategory closed under direct sums and containing \(D^b_{\mathrm{ctf}}(X,\Lambda)\), and hence they form all of \(\Loc(D^b_{\mathrm{ctf}}(X,\Lambda))\).

Extension by zero is exact and left adjoint to \(j^*\), so it preserves direct sums. It preserves constructibility, since we may add the closed complement as a stratum on which the sheaf is zero. For finite \(A\) it also preserves finite Tor amplitude, as one checks on stalks, while for \(\mathcal O_E\), \(E\) and \(\overline{\mathbf Q}_\ell\) we have \(D^b_{\mathrm{ctf}}=\Dbc\), so constructibility suffices; for \(\mathcal O_E\) this is Lemma \hyperref[bp68]{5.20}. With this, the localizing-subcategory argument shows that \(j_!\) restricts as stated.
\end{proof}

% text: prose-05-28
\phantomsection\label{bp78}%
The proper base change map and the comparison maps between compactifications are natural transformations between exact functors on \(\Loc(D^b_{\mathrm{ctf}}(-,\Lambda))\) that preserve direct sums. To show that such a transformation is an isomorphism, it suffices to check this on \(D^b_{\mathrm{ctf}}(-,\Lambda)\).

% object: lemma-checking-isomorphisms-constructible-complexes (blueprint/chap-Construction-of-f/sec-Direct-image-with-compact-support/subsec-Construction-of-Rf/lemma-checking-isomorphisms-constructible-complexes/_section.md)
\phantomsection\label{bp79}%
\begin{bpstatementit}{Lemma~5.23}{ (Checking isomorphisms on constructible complexes)}
Let \(\Lambda\) be admissible, let \(F_1,F_2:\Loc(D^b_{\mathrm{ctf}}(X,\Lambda))\to\Loc(D^b_{\mathrm{ctf}}(Y,\Lambda))\) be exact functors preserving direct sums, and let \(\eta:F_1\to F_2\) be a natural transformation of exact functors. If \(\eta_K\) is an isomorphism for every \(K\in D^b_{\mathrm{ctf}}(X,\Lambda)\), then it is an isomorphism for every \(K\in\Loc(D^b_{\mathrm{ctf}}(X,\Lambda))\).
\end{bpstatementit}

% proof: proof/lemma-checking-isomorphisms-constructible-complexes
\begin{proof}[{Proof}]\phantomsection\label{bp80}%
The objects \(K\) for which \(\eta_K\) is an isomorphism are closed under shifts and cones by the five lemma, and under direct sums because \(F_1\) and \(F_2\) preserve them. Since they contain \(D^b_{\mathrm{ctf}}(X,\Lambda)\), they form all of \(\Loc(D^b_{\mathrm{ctf}}(X,\Lambda))\).
\end{proof}

% text: prose-05-29
\phantomsection\label{bp81}%
For a cartesian square with \(p\) proper, the base change map \(g^*Rp_*\to Rp'_*g'^*\) of (4.5) is a map of functors \(D(X_{\mathrm{pro}\acute{e}t},\Lambda)\to D(Y'_{\mathrm{pro}\acute{e}t},\Lambda)\). The lemma lets us pass from \(D^b_{\mathrm{ctf}}(X,\Lambda)\), where it is an isomorphism, to \(\Loc(D^b_{\mathrm{ctf}}(X,\Lambda))\).

% object: proposition-proper-base-change-localizing-categories (blueprint/chap-Construction-of-f/sec-Direct-image-with-compact-support/subsec-Construction-of-Rf/proposition-proper-base-change-localizing-categories/_section.md)
\phantomsection\label{bp82}%
\begin{bpstatementit}{Proposition~5.24}{ (Proper base change on the localizing categories)}
Let \(\Lambda\) be admissible. Consider a cartesian square

\[
\begin{array}{ccc}
X'&\xrightarrow{g'}&X\\
\downarrow p'&&\downarrow p\\
Y'&\xrightarrow{g}&Y
\end{array}
\]

with \(p\) proper. The base change map \(g^*Rp_*K\to Rp'_*g'^*K\) is an isomorphism for every \(K\in\Loc(D^b_{\mathrm{ctf}}(X,\Lambda))\).
\end{bpstatementit}

% proof: proof/proposition-proper-base-change-localizing-categories
\begin{proof}[{Proof}]\phantomsection\label{bp83}%
On \(D^b_{\mathrm{ctf}}(X,\Lambda)\) the base change map is an isomorphism by Proposition \bpout{4.40}, so by Lemma \hyperref[bp79]{5.23} it suffices to show that \(g^*Rp_*\) and \(Rp'_*g'^*\) are exact functors \(\Loc(D^b_{\mathrm{ctf}}(X,\Lambda))\to\Loc(D^b_{\mathrm{ctf}}(Y',\Lambda))\) preserving direct sums. For \(Rp_*\) and \(Rp'_*\) this is Proposition \hyperref[bp76]{5.22}, so it remains to treat inverse image. It is exact and preserves direct sums, and it sends \(D^b_{\mathrm{ctf}}(-,\Lambda)\) into \(D^b_{\mathrm{ctf}}(-,\Lambda)\). Indeed, for finite \(A\) this holds because of bounded flat constructible representatives, and for \(\mathcal O_E\), \(E\) and \(\overline{\mathbf Q}_\ell\) it holds because pullback preserves the local description of constructible sheaves on strata (Lemma \hyperref[bp68]{5.20} in the integral case). Given this, the localizing-subcategory argument before Proposition \hyperref[bp76]{5.22} shows that inverse image also preserves \(\Loc(D^b_{\mathrm{ctf}}(-,\Lambda))\).
\end{proof}

% object: remark-compactifications (blueprint/chap-Construction-of-f/sec-Direct-image-with-compact-support/subsec-Construction-of-Rf/remark-compactifications/_section.md)
\phantomsection\label{bp84}%
\begin{bpstatement}{Remark~5.25}{ (Compactifications)}
\begin{enumerate}
\item By Nagata's theorem (Theorem \bpout{A.2}), every separated morphism of finite type \(f:X\to Y\) with \(Y\) qcqs has a compactification, that is, a factorization \(f=p\circ j\) with \(j:X\to\overline X\) a quasi-compact open immersion and \(p:\overline X\to Y\) proper.
\item In comparisons we may replace \(\overline X\) by the scheme-theoretic closure of \(X\), so that \(X\) is dense, and by Lemma \bpout{A.5} this does not change \(Rp_*j_!\).
\end{enumerate}
\end{bpstatement}

% text: prose-05-30
\phantomsection\label{bp85}%
The comparisons between compactifications below are the proper direct images of the unique maps out of \(j_!K\) that restrict to the identity on \(X\). They are unique because, for an open immersion \(j:X\to Z\), the functor \(j_!\) is left adjoint to \(j^*\) on the pro-étale derived categories, so that a map out of \(j_!K\) is determined by its restriction to \(X\).

% object: proposition-construction-independence-rf-shriek (blueprint/chap-Construction-of-f/sec-Direct-image-with-compact-support/subsec-Construction-of-Rf/proposition-construction-independence-rf-shriek/_section.md)
\phantomsection\label{bp86}%
\begin{bpstatementit}{Proposition~5.26}{ (Construction and independence of $Rf_!$)}
Let \(\Lambda\) be admissible and let \(f:X\to Y\) be separated of finite type, with compactification \(f=p\circ j\) as above. The functor

\[
Rf_!=Rp_*\circ j_!:\Loc(D^b_{\mathrm{ctf}}(X,\Lambda))\to\Loc(D^b_{\mathrm{ctf}}(Y,\Lambda))
\]

is exact, preserves direct sums, and maps \(D^b_{\mathrm{ctf}}(X,\Lambda)\) into \(D^b_{\mathrm{ctf}}(Y,\Lambda)\). Two compactifications give canonically isomorphic functors, and these isomorphisms satisfy the cocycle condition.
\end{bpstatementit}

% proof: proof/proposition-construction-independence-rf-shriek
\begin{proof}[{Proof}]\phantomsection\label{bp87}%
For a single compactification there is nothing new to prove, and two compactifications are compared through a third one that refines both, by the argument described at the start of this subsection.

The properties of \(Rp_*\circ j_!\) for a single compactification follow from Proposition \hyperref[bp76]{5.22} and the preservation of \(D^b_{\mathrm{ctf}}(-,\Lambda)\) by proper direct image (Proposition \bpout{4.40}).

For independence it suffices to compare a compactification with a refinement of it. Indeed, the schematic closure of \(X\) in a compactification is a refinement of it, since its inverse image over \(X\) is \(X\). Moreover, any two dense compactifications have a common refinement (Proposition \bpout{A.7}).

\emph{Refinements.} Let \(j:X\to Z\) and \(j':X\to Z'\) be compactifications over \(Y\), and let \(h:Z'\to Z\) be a refinement (Definition \bpout{A.1}), that is, a proper \(Y\)-morphism with \(h\circ j'=j\) and \(h^{-1}(X)=X\). To construct the comparison map, note that restriction to the open subscheme \(X\) commutes with \(Rh_*\). This holds because \(h^{-1}(X)=X\) and restriction is exact and preserves K-injectives, its left adjoint \(j_!\) being exact. As \(j'^*j'_!K=K\), we get \(j^*Rh_*j'_!K\simeq K\). By adjunction this isomorphism corresponds to a map

\[
e_{h,K}:j_!K\longrightarrow Rh_*j'_!K
\tag{5.7}
\]

whose restriction to \(X\) is the identity. We claim that it is an isomorphism for \(K\in\Loc(D^b_{\mathrm{ctf}}(X,\Lambda))\). By Proposition \hyperref[bp76]{5.22}, both sides are exact functors of \(K\) preserving direct sums, so by Lemma \hyperref[bp79]{5.23} we may assume that \(K\in D^b_{\mathrm{ctf}}(X,\Lambda)\), where proper base change is available.

By the exact sequence for an open and its closed complement, applied in each degree [Stacks, \href{https://stacks.math.columbia.edu/tag/09AH}{09AH}], a morphism in \(D(Z_{\mathrm{pro}\acute{e}t},\Lambda)\) that is an isomorphism on \(X\) and on \(Z\setminus X\) is an isomorphism. On \(X\) our map is the identity, so it remains to see that on \(Z\setminus X\) both sides vanish. For \(j_!K\) this holds by definition of extension by zero, while for \(Rh_*j'_!K\) we use proper base change on \(D^b_{\mathrm{ctf}}\) (Proposition \hyperref[bp82]{5.24}) for the inclusion of \(Z\setminus X\) into \(Z\). It shows that the restriction of \(Rh_*j'_!K\) is the direct image of the restriction of \(j'_!K\) to \(h^{-1}(Z\setminus X)=Z'\setminus X\), which is zero. This proves the claim. Applying \(Rp_*\) and using \(R(p\circ h)_*=Rp_*Rh_*\), we obtain an isomorphism between the two candidate functors \(Rf_!\), namely \(Rp_*j_!\) and \(R(ph)_*j'_!\).

If \(Z''\xrightarrow{k}Z'\xrightarrow{h}Z\) are two successive refinements, with \(j'':X\to Z''\) the open immersion, then \(Rh_*(e_{k,K})\circ e_{h,K}\) and \(e_{hk,K}\) are maps from \(j_!K\) to \(R(hk)_*j''_!K\) that both restrict to the identity on \(X\). By the observation before the proposition, they are equal.

\emph{Common refinements.} Two arbitrary compactifications need not be related by a refinement, but they have a common one. Indeed, given dense compactifications \(Z_1\) and \(Z_2\), Proposition \bpout{A.7} provides a dense compactification \(Z\) with proper maps \(h_i:Z\to Z_i\) satisfying \(h_i^{-1}(X)=X\), and similarly for finitely many compactifications. The refinement isomorphisms for \(h_1\) and \(h_2\) then compose to an isomorphism between the functors defined by \(Z_1\) and \(Z_2\), and transitivity of the refinement maps on a common refinement shows that this comparison does not depend on the chosen refinement and satisfies the cocycle condition.
\end{proof}

% object: example-cuspidal-compactification (blueprint/chap-Construction-of-f/sec-Direct-image-with-compact-support/subsec-Construction-of-Rf/example-cuspidal-compactification/_section.md)
\phantomsection\label{bp88}%
\begin{bpstatement}{Example~5.27}{ (A singular compactification of the line)}
Over \(k\), let \(C\subset\mathbf P^2\) be the cuspidal cubic \(zy^2=x^3\), with cusp \(c=(0:0:1)\). Its normalization is the map \(h:\mathbf P^1\to C\), \((s:u)\mapsto(su^2:u^3:s^3)\). It sends \(u=0\) to the cusp and the affine line \(u\neq0\) isomorphically onto \(C\setminus\{c\}\). So \(X=\mathbf A^1\) has the two compactifications \(\mathbf P^1\) and \(C\). Since \(h^{-1}(C\setminus\{c\})=\mathbf A^1\), the map \(h\) is a refinement. The refinement map \(j_!\Lambda\to Rh_*j'_!\Lambda\) is the identity on \(X\), and at the cusp both stalks vanish: for the right side, proper base change shows that the stalk of \(Rh_*j'_!\Lambda\) at \(c\) is the cohomology of \(j'_!\Lambda\) on the single point \(h^{-1}(c)=\{\infty\}\), where \(j'_!\Lambda\) is zero. So the singularity of \(C\) does not affect \(R\Gamma_c(\mathbf A^1,\Lambda)\).
\end{bpstatement}

% text: prose-05-31
\phantomsection\label{bp89}%
To make \(Rf_!\) functorial in \(f\), and for the base change theorems later on, we need composition and base change isomorphisms. The base change map is defined as follows. Let \(f:X\to Y\) be separated of finite type with compactification \(f=p\circ j\), let \(v:Y'\to Y\) be a morphism, and let \(f':X'\to Y'\), \(u:X'\to X\), \(p':\overline X'\to Y'\), \(j':X'\to\overline X'\) and \(v':\overline X'\to\overline X\) be the base changes, so that \(f'=p'\circ j'\) is a compactification. For \(K\in D(X_{\mathrm{pro}\acute{e}t},\Lambda)\), the isomorphism \(j'^*v'^*j_!K\simeq u^*j^*j_!K=u^*K\) corresponds under \(j'_!\dashv j'^*\) to a map \(j'_!u^*K\to v'^*j_!K\). This map is an isomorphism: it is the identity on the open \(X'\), and both sides vanish on the closed complement \(v'^{-1}(\overline X\setminus X)=\overline X'\setminus X'\) [Stacks, \href{https://stacks.math.columbia.edu/tag/09AH}{09AH}]. The base change map for \(Rf_!\) is the composite

\[
v^*Rp_*j_!K\longrightarrow Rp'_*v'^*j_!K\simeq Rp'_*j'_!u^*K
\tag{5.8}
\]

of the base change map (4.5) for \(p\) and \(Rp'_*\) of the inverse of this isomorphism.

% object: proposition-composition-base-change-rf-shriek (blueprint/chap-Construction-of-f/sec-Direct-image-with-compact-support/subsec-Construction-of-Rf/proposition-composition-base-change-rf-shriek/_section.md)
\phantomsection\label{bp90}%
\begin{bpstatementit}{Proposition~5.28}{ (Composition and base change for $Rf_!$)}
For every admissible \(\Lambda\), the following statements hold on \(\Loc(D^b_{\mathrm{ctf}}(-,\Lambda))\) and restrict to \(D^b_{\mathrm{ctf}}(-,\Lambda)\).

\begin{enumerate}
\item For separated morphisms of finite type \(X\xrightarrow fY
  \xrightarrow gZ\) there is a canonical isomorphism \(Rg_!\circ Rf_!\simeq R(gf)_!\). These isomorphisms are associative and compatible with the identity \(R(\mathrm{id})_!=\mathrm{id}\).
\item Let \(v:Y'\to Y\) be a morphism, \(f':X'\to Y'\) the base change of \(f\) along \(v\), and \(u:X'\to X\) the projection. The base change map (5.8) is an isomorphism \(v^*Rf_!\simeq Rf'_!u^*\). It does not depend on the compactification, and it is compatible with composing cartesian squares and with the composition isomorphisms of (1).
\end{enumerate}
\end{bpstatementit}

% proof: proof/proposition-composition-base-change-rf-shriek
\begin{proof}[{Proof}]\phantomsection\label{bp91}%
For composition, the only new ingredient is a map that moves an extension by zero past a proper direct image. It is an isomorphism by the same argument as for refinements: it restricts to the identity on an open, and both sides vanish on the closed complement by proper base change. For base change of \(Rf_!\) we combine proper base change with the elementary fact that extension by zero commutes with pullback.

\emph{The composition isomorphism.} Choose compactifications \(f=p\circ j\) and \(g=q\circ j_Y\), with \(j:X\to\overline X\) and \(j_Y:Y\to\overline Y\) open and \(p,q\) proper. In the composite \(Rq_*j_{Y!}Rp_*j_!\), the middle factor \(j_{Y!}Rp_*\) applies a proper direct image before an extension by zero, and we want to interchange the two. To do this, compactify the morphism \(j_Y\circ p:\overline X\to\overline Y\) as \(\overline X\xrightarrow{j'}W\xrightarrow{r}\overline Y\), with \(j'\) open, \(r\) proper and \(\overline X\) dense in \(W\). Then \(r^{-1}(Y)=\overline X\), because the open immersion \(\overline X\to r^{-1}(Y)\) is a morphism from a proper \(Y\)-scheme to a separated \(Y\)-scheme, hence proper, and its image is closed, open and dense.

The interchange map is defined as follows. For \(M\) in \(D(\overline X_{\mathrm{pro}\acute{e}t},\Lambda)\), restriction to the open \(Y\subset\overline Y\) identifies \(j_Y^*Rr_*j'_!M\) with \(Rp_*M\), because \(r^{-1}(Y)=\overline X\), the restriction of \(r\) over \(Y\) is \(p\), and \(j'_!M\) restricts to \(M\) on \(\overline X\). By the adjunction \(j_{Y!}\dashv j_Y^*\), we obtain a map

\[
\iota_M:j_{Y!}Rp_*M\longrightarrow Rr_*j'_!M,
\tag{5.9}
\]

which is the unique map restricting to the identity of \(Rp_*M\) over \(Y\).

For \(M\in D^b_{\mathrm{ctf}}(\overline X,\Lambda)\), this map is an isomorphism. Indeed, over \(Y\) it is the identity, and over \(\overline Y\setminus Y\) both sides vanish. For the right-hand side, the vanishing follows from proper base change on \(D^b_{\mathrm{ctf}}\) (Proposition \hyperref[bp82]{5.24}) for the inclusion of \(\overline Y\setminus Y\) into \(\overline Y\), since \(j'_!M\) vanishes on \(r^{-1}(\overline Y\setminus Y)=W\setminus\overline X\). Both sides preserve direct sums on \(\Loc(D^b_{\mathrm{ctf}}(\overline X,\Lambda))\) (Proposition \hyperref[bp76]{5.22}). Hence Lemma \hyperref[bp79]{5.23} shows that \(\iota_M\) is an isomorphism for all \(M\in\Loc(D^b_{\mathrm{ctf}}(\overline X,\Lambda))\). For \(M=j_!K\) it gives

\[
Rq_*j_{Y!}Rp_*j_!K\simeq Rq_*Rr_*j'_!j_!K=R(qr)_*(j'j)_!K,
\]

and the right side is \(R(gf)_!K\) computed with the compactification \(gf=(qr)\circ(j'j)\).

\emph{Verification of canonicity and associativity.} Every map used is characterized by its restriction to an open subscheme, since the refinement maps (5.7) and the maps \(\iota_M\) restrict to identities. For example, replacing any of the chosen compactifications by a refinement commutes with \(\iota_M\), because both composites restrict to the same identity and are then equal by adjunction. For associativity, given three composable morphisms, we choose common refinements (by graph closure) of the finitely many compactifications that occur in the composites of two of them. On these, each of the two bracketings is built from composites of proper direct images, composites of extensions by zero, and maps of the form \(\iota_M\). The first two kinds of pieces are associative: the isomorphisms \(Rq_*Rr_*\simeq R(qr)_*\) are associative because they are the isomorphisms of derived functors of composites of left exact functors (equivalently, they are adjoint to the composition isomorphisms of the inverse images), and composition of extensions by zero is associative by adjunction. Moreover, interchanging two adjacent open and proper steps in either order gives the same map by the uniqueness of \(\iota_M\). So the two bracketings agree. Similarly, the unit conditions are the special case where one compactification is the identity. Finally, transitivity of the refinement maps carries these equalities over to the originally chosen compactifications.

\emph{Base change.} The base change map for \(Rf_!=Rp_*j_!\) is built from one for each factor, so we treat \(j_!\) and \(Rp_*\) separately. Since open immersions and proper morphisms are stable under base change, we may pull back a compactification \(f=p\circ j\) along \(v\) to a compactification \(f'=p'\circ j'\), where now \(j'\) and \(p'\) denote the base changes of \(j\) and \(p\). For \(j_!\), extension by zero commutes with inverse image, as shown before the proposition. Combining this with proper base change on \(\Loc(D^b_{\mathrm{ctf}}(\overline X,\Lambda))\) (Proposition \hyperref[bp82]{5.24}), we obtain the isomorphism \(v^*Rp_*j_!\simeq Rp'_*v'^*j_!\simeq Rp'_*j'_!u^*\) on \(\Loc(D^b_{\mathrm{ctf}}(X,\Lambda))\). This isomorphism is compatible with change of compactification because the refinement maps are characterized by their restriction to the open subscheme. It is also compatible with composing two cartesian squares, because both factors are. For the first factor, a map between extensions by zero from an open is determined by its restriction to that open, since \(\Hom(j_!B,j_!B')=\Hom(B,B')\). For the second, a base change map is the adjoint transpose of a counit, and the counit of a composite adjunction is the composite of the counits. Finally, compatibility with the composition isomorphisms of (1) follows in the same way as canonicity: both composites are built from refinement maps, maps \(\iota_M\) and proper base change maps, and they agree because their restrictions to the open part are the same identity.
\end{proof}

% text: prose-05-32
\phantomsection\label{bp92}%
On a space, a section with compact support is in particular a section, and so compactly supported cohomology maps to ordinary cohomology. The algebraic counterpart is a map \(Rf_!\to Rf_*\), which is an isomorphism for proper \(f\), and we need it to be compatible with the composition maps just constructed: in Section \bpout{6.1} this is one of the compatibilities that the comparison with the étale site has to respect.

% object: proposition-forgetting-supports-composition (blueprint/chap-Construction-of-f/sec-Direct-image-with-compact-support/subsec-Construction-of-Rf/proposition-forgetting-supports-composition/_section.md)
\phantomsection\label{bp93}%
\begin{bpstatementit}{Proposition~5.29}{ (Forgetting supports and composition)}
There is a natural map \(\eta_f:Rf_!\to Rf_*\) on \(\Loc(D^b_{\mathrm{ctf}}(X,\Lambda))\). For a compactification \(f=pj\) it is

\[
Rp_*j_!\longrightarrow Rp_*Rj_*=Rf_*,
\]

where \(j_!\to Rj_*\) is adjoint to the inverse of the counit \(j^*Rj_*\xrightarrow{\sim}\mathrm{id}\). It is independent of the compactification, and the square

\[
\begin{CD}
Rg_!Rf_! @>{c^!_{g,f}}>> R(gf)_!\\
@V{Rg_*(\eta_f)\,\circ\,(\eta_g Rf_!)}VV @VV{\eta_{gf}}V\\
Rg_*Rf_* @>{c^*_{g,f}}>> R(gf)_*
\end{CD}
\]

commutes. Here \(c^!_{g,f}\) is the composition map of Proposition \hyperref[bp90]{5.28}, and \(c^*_{g,f}\) is the composition map for ordinary direct images. If \(f\) is proper, \(\eta_f\) is an isomorphism; under these isomorphisms the two composition maps agree whenever \(f\) and \(g\) are proper.
\end{bpstatementit}

% proof: proof/proposition-forgetting-supports-composition
\begin{proof}[{Proof}]\phantomsection\label{bp94}%
All maps involved are maps out of an extension by zero, so each of them is determined by its restriction to the open subscheme, where it is an identity or an inverse counit. We use this repeatedly.

Independence of the compactification reduces to a refinement \(h:Z'\to Z\). The two maps \(j_!K\to Rh_*Rj'_*K=Rj_*K\) obtained from (5.7) and from \(j_!K\to Rj_*K\) have the same adjoint under \(j_!\dashv j^*\), namely the inverse of \(j^*Rj_*K\to K\). Hence they are equal. Applying \(Rp_*\) gives independence.

The square is checked on the compactifications that define \(c^!_{g,f}\) and the interchange map (5.9). Apart from composition and naturality squares for ordinary direct images, it comes down to comparing the two maps

\[
j_{Y!}Rp_*M\longrightarrow Rr_*j'_!M
\longrightarrow Rr_*Rj'_*M=Rj_{Y*}Rp_*M,
\qquad
j_{Y!}Rp_*M\longrightarrow Rj_{Y*}Rp_*M.
\]

Their adjoints under \(j_{Y!}\dashv j_Y^*\) are both the inverse counit for \(j_Y^*\dashv Rj_{Y*}\), because the interchange map restricts to the identity on \(Y\). So the two maps agree. Applying \(Rq_*\) with \(M=j_!K\), and using naturality for \(j_!K\to Rj_*K\), we get the displayed square for the chosen compactifications. The refinement calculation extends it to any choices. If \(f\) is proper, the compactification with \(j=\mathrm{id}\) has \(\eta_f\) equal to the identity, and the last assertion follows from the square.
\end{proof}



% object: sec-Brown-representability (blueprint/chap-Construction-of-f/sec-Brown-representability/_section.md), content hidden
\subsection*{\texorpdfstring{5.5\quad Brown representability}{5.5 Brown representability}}
\phantomsection\label{bp95}%
\addcontentsline{toc}{subsection}{\texorpdfstring{5.5\quad Brown representability}{5.5 Brown representability}}
% text: prose-05-33
\phantomsection\label{bp96}%
Let \(\Lambda\) be an admissible coefficient ring and fix \(N\) in \(\Loc(D^b_{\mathrm{ctf}}(Y,\Lambda))\). Since the functor \(Rf_!\) preserves direct sums (Proposition \hyperref[bp86]{5.26}), the functor \(M\mapsto\Hom(Rf_!M,N)\) turns direct sums into products. If the objects of \(D^b_{\mathrm{ctf}}(X,\Lambda)\) are compact generators of the localizing category, Brown representability produces an object \(f^!N\) representing it. The next lemma checks this hypothesis for every admissible \(\Lambda\); for \(\mathcal O_E\) it uses the compactness proved in Subsection \hyperref[bp63]{5.4.2} by reduction to torsion coefficients.

Compactness enters a second time. The functor \(Rf_!\) also maps \(D^b_{\mathrm{ctf}}(X,\Lambda)\) into \(D^b_{\mathrm{ctf}}(Y,\Lambda)\), so it sends the compact generators to compact objects. For this reason the adjoint \(f^!\) preserves direct sums (Theorem \hyperref[bp100]{5.31}).

% object: lemma-hypotheses-brown-representability (blueprint/chap-Construction-of-f/sec-Brown-representability/lemma-hypotheses-brown-representability/_section.md)
\phantomsection\label{bp97}%
\begin{bpstatementit}{Lemma~5.30}{ (Hypotheses of Brown representability)}
For every admissible \(\Lambda\), the category \(\Loc(D^b_{\mathrm{ctf}}(X,\Lambda))\) is locally small and triangulated, has arbitrary direct sums, and has split idempotents. The thick category \(D^b_{\mathrm{ctf}}(X,\Lambda)\) is essentially small, and a small skeleton is a set of compact generators of \(\Loc(D^b_{\mathrm{ctf}}(X,\Lambda))\).
\end{bpstatementit}

% proof: proof/lemma-hypotheses-brown-representability
\begin{proof}[{Proof}]\phantomsection\label{bp98}%
The set-theoretic conditions are general facts about derived categories of Grothendieck abelian categories. Compactness, on the other hand, was proved separately for each kind of coefficients, and generation follows from the definition of the localizing category.

As the derived category of a Grothendieck abelian category of module sheaves, the ambient category is locally small, and so is its full subcategory \(\Loc(D^b_{\mathrm{ctf}}(-,\Lambda))\). Moreover, the latter has direct sums by definition, so its idempotents split [Stacks, \href{https://stacks.math.columbia.edu/tag/05QW}{05QW}]. Finally, the subcategory \(D^b_{\mathrm{ctf}}(X,\Lambda)\) is thick, since constructibility and finite Tor amplitude are preserved by shifts, cones and direct summands.

To see that \(D^b_{\mathrm{ctf}}(X,\Lambda)\) is essentially small, note first that the constructible sheaves in the heart have a set of representatives for each admissible \(\Lambda\). A bounded complex with such cohomology is built from its cohomology sheaves by finitely many extensions, formed with truncations in the ambient category. Each extension step is classified by an element of a Hom set of the ambient category, and these Hom sets are sets because the ambient category is locally small. Hence these complexes form an essentially small category, of which \(D^b_{\mathrm{ctf}}(X,\Lambda)\) is a full subcategory.

Compactness for finite \(A\) is Proposition \hyperref[bp45]{5.14}. For \(E\) and \(\overline{\mathbf Q}_\ell\) it is Theorem \hyperref[bp40]{5.13}, and for \(\mathcal O_E\) it is Proposition \hyperref[bp71]{5.21}; in the integral and rational cases, \(D^b_{\mathrm{ctf}}=\Dbc\).

For generation, suppose that \(K\in\Loc(D^b_{\mathrm{ctf}}(X,\Lambda))\) receives no nonzero maps from shifts of objects of a small skeleton. The objects \(M\) with \(\Hom(M[n],K)=0\) for all \(n\) form a localizing subcategory containing \(D^b_{\mathrm{ctf}}\). Hence this subcategory contains \(K\). Then \(\mathrm{id}_K=0\), and so \(K=0\).
\end{proof}

% text: prose-05-34
\phantomsection\label{bp99}%
With the hypotheses checked, the theorem below follows from Brown representability in the form recalled at the start of the chapter.

% object: theorem-existence-exceptional-inverse-image (blueprint/chap-Construction-of-f/sec-Brown-representability/theorem-existence-exceptional-inverse-image/_section.md)
\phantomsection\label{bp100}%
\begin{bpstatementit}{Theorem~5.31}{ (Existence of the exceptional inverse image)}
Let \(\Lambda\in\{A,\mathbf Z_\ell,\mathcal O_E,\mathbf Q_\ell,E,\overline{\mathbf Q}_\ell\}\) be any admissible coefficient ring. For every separated morphism of finite type \(f:X\to Y\), the coproduct-preserving functor \(Rf_!:\Loc(D^b_{\mathrm{ctf}}(X,\Lambda))\to\Loc(D^b_{\mathrm{ctf}}(Y,\Lambda))\) has an exact right adjoint

\[
f^!:\Loc(D^b_{\mathrm{ctf}}(Y,\Lambda))\longrightarrow\Loc(D^b_{\mathrm{ctf}}(X,\Lambda)),
\qquad
\Hom(Rf_!M,N)\simeq\Hom(M,f^!N),
\]

for \(M\in\Loc(D^b_{\mathrm{ctf}}(X,\Lambda))\) and \(N\in\Loc(D^b_{\mathrm{ctf}}(Y,\Lambda))\). It is unique up to unique isomorphism compatible with the adjunction, and it preserves arbitrary direct sums.
\end{bpstatementit}

% proof: proof/theorem-existence-exceptional-inverse-image
\begin{proof}[{Proof}]\phantomsection\label{bp101}%
We construct \(f^!N\) one object at a time as the object representing maps out of \(Rf_!(-)\) into \(N\), and then check direct sums on the compact generators.

For \(N\in\Loc(D^b_{\mathrm{ctf}}(Y,\Lambda))\), the functor

\[
\Loc(D^b_{\mathrm{ctf}}(X,\Lambda))^{\mathrm{op}}\longrightarrow\operatorname{Ab},
\qquad M\longmapsto\Hom(Rf_!M,N)
\]

is cohomological and sends direct sums to products. By Lemma \hyperref[bp97]{5.30}, \(\Loc(D^b_{\mathrm{ctf}}(X,\Lambda))\) is compactly generated, and Brown representability [Stacks, \href{https://stacks.math.columbia.edu/tag/0A8F}{0A8F}] says that on such a category every cohomological functor sending direct sums to products is representable, say by \(f^!N\). By the Yoneda lemma, these objects form a right adjoint, and a right adjoint of an exact functor of triangulated categories is exact [Stacks, \href{https://stacks.math.columbia.edu/tag/0A8G}{0A8G}].

To see that \(f^!\) preserves direct sums, we test the canonical map \(\bigoplus_\alpha f^!N_\alpha\to f^!(\bigoplus_\alpha N_\alpha)\) on a compact generator \(G\in D^b_{\mathrm{ctf}}(X,\Lambda)\) and its shifts. On the source, compactness of \(G\) and adjunction give

\[
\Hom\Bigl(G[n],\bigoplus_\alpha f^!N_\alpha\Bigr)
=\bigoplus_\alpha\Hom(G[n],f^!N_\alpha)
=\bigoplus_\alpha\Hom(Rf_!G[n],N_\alpha).
\]

On the target, adjunction and compactness of \(Rf_!G\in D^b_{\mathrm{ctf}}(Y,\Lambda)\) give the same group. The map therefore induces bijections on \(\Hom(G[n],-)\) for all compact generators \(G\) and all \(n\). Applying the detection property to its cone, we see that it is an isomorphism.
\end{proof}


% object: sec-Composition-and-restriction (blueprint/chap-Construction-of-f/sec-Composition-and-restriction/_section.md), content hidden
\subsection*{\texorpdfstring{5.6\quad Composition and restriction}{5.6 Composition and restriction}}
\phantomsection\label{bp102}%
\addcontentsline{toc}{subsection}{\texorpdfstring{5.6\quad Composition and restriction}{5.6 Composition and restriction}}
% text: prose-05-35
\phantomsection\label{bp103}%
Brown representability gives \(f^!\) one morphism at a time. To compute it in Chapter \bpout{6}, however, we factor \(f\) locally. The two are linked by uniqueness of adjoints: an isomorphism between composites of left adjoints, such as \(Rg_!\circ Rf_!\simeq R(gf)_!\) or \(Rf_!=Rp_*\circ j_!\), gives an isomorphism between the right adjoints, composed in the opposite order. Since \(j^!=j^*\) for an open immersion \(j\), it follows that \(f^!\) is local on \(X\) and on \(Y\). For example, a smooth \(f\) is locally the composite of an étale map and a projection \(\mathbf A^n_Y\to Y\). By composition and locality, \(f^!\) is then determined by its values on these two kinds of maps.

Base change for \(f^!\) does not follow in this way. Passing to right adjoints in the base change isomorphism \(v^*Rf_!\simeq Rf'_!u^*\) relates \(f^!Rv_*\) and \(Ru_*f'^!\). These involve direct images, not the inverse images \(u^*f^!\) and \(f'^!v^*\) that one wants to compare. We prove base change for smooth morphisms in Proposition \bpout{6.9}. That \(f^!\) preserves \(D^b_{\mathrm{ctf}}\) is proved in Theorems \bpout{6.20} and \bpout{Theorem 6.27}.

% object: proposition-composition-open-restriction-compactifications-f (blueprint/chap-Construction-of-f/sec-Composition-and-restriction/proposition-composition-open-restriction-compactifications-f/_section.md)
\phantomsection\label{bp104}%
\begin{bpstatementit}{Proposition~5.32}{ (Composition, open restriction and compactifications for $f^!$)}
For every admissible \(\Lambda\), the adjoints \(f^!\) on \(\Loc(D^b_{\mathrm{ctf}}(-,\Lambda))\) have the following canonical compatibilities.

\begin{enumerate}
\item For composable \(X\xrightarrow fY\xrightarrow gZ\) there are canonical isomorphisms \((gf)^!\simeq f^!g^!\), associative and compatible with the identity.
\item For an open immersion \(j:U\to X\), \(j^!\simeq j^*\).
\item Let \(v:V\to Y\) be an open immersion, \(X_V=f^{-1}(V)\) with inclusion \(u:X_V\to X\), and \(f_V:X_V\to V\) the restriction of \(f\). Then \(u^*f^!\simeq f_V^!v^*\).
\item If \(f=p\circ j\) is a compactification, then \(f^!\simeq j^*\circ p^!\) on \(\Loc(D^b_{\mathrm{ctf}}(Y,\Lambda))\).
\end{enumerate}

In particular, \(f^!\) can be computed locally on \(X\) and on \(Y\) for open coverings.
\end{bpstatementit}

% proof: proof/proposition-composition-open-restriction-compactifications-f
\begin{proof}[{Proof}]\phantomsection\label{bp105}%
Each statement is the right adjoint of a statement about lower shriek functors already proved, and we conclude by uniqueness of adjoints.

(1) For \(M\in\Loc(D^b_{\mathrm{ctf}}(X,\Lambda))\) and \(N\in\Loc(D^b_{\mathrm{ctf}}(Z,\Lambda))\),

\[
\Hom(M,f^!g^!N)
\simeq\Hom(Rf_!M,g^!N)
\simeq\Hom(Rg_!Rf_!M,N)
\simeq\Hom(R(gf)_!M,N),
\]

where the last step holds by Proposition \hyperref[bp90]{5.28}. By the Yoneda lemma, \(f^!g^!\simeq(gf)^!\). Passing to right adjoints transforms the associativity and identity constraints of Proposition proposition-composition-base-change-rf-shriek into those for \(f^!\).

(2) For an open immersion, \(Rj_!=j_!\) (take the trivial compactification \(p=\mathrm{id}\)), and \(j_!\dashv j^*\) on the localizing categories; here we use that \(j^*\) preserves the localizing categories, as inverse images do (see the proof of Proposition \hyperref[bp82]{5.24}). By uniqueness of right adjoints, \(j^!\simeq j^*\).

(3) Since \(f\circ u=v\circ f_V\), parts (1) and (2) show that \(u^*f^!\simeq(fu)^!=(vf_V)^!\simeq f_V^!v^*\). For an open \(j:U\to X\), parts (1) and (2) also give \(j^*f^!\simeq(fj)^!\). These two isomorphisms show that \(f^!\) is local on \(X\) and on \(Y\).

(4) On \(\Loc(D^b_{\mathrm{ctf}}(-,\Lambda))\), \(j^*\) is right adjoint to \(j_!\) and \(p^!\) to \(Rp_*=Rp_!\), so \(j^*\circ p^!\) and \(f^!\) are both right adjoint to \(Rp_*\circ j_!=Rf_!\), and they are isomorphic by uniqueness of right adjoints.
\end{proof}

% text: prose-05-36
\phantomsection\label{bp106}%
Once we know that \(f^!\) preserves \(D^b_{\mathrm{ctf}}(-,\Lambda)\) (Chapter \bpout{6}), these isomorphisms restrict to the full subcategories \(D^b_{\mathrm{ctf}}(-,\Lambda)\), and the same statements hold for the adjunction between \(D^b_{\mathrm{ctf}}(X,\Lambda)\) and \(D^b_{\mathrm{ctf}}(Y,\Lambda)\).



\clearpage\phantomsection\hypertarget{bp-outside}{}%
\section*{Outside this export}
\addcontentsline{toc}{section}{Outside this export}
References marked \ensuremath{{}^\ast} are to parts of the blueprint this document leaves out:
\begin{description}
\item[Section 1.9] Conventions, in 1 Introduction (\bpslug{sec-{}\allowbreak{}Conventions})
\item[Theorem 2.17] Enough affine weakly contractible objects, in Ⅰ The six operations (\bpslug{theorem-{}\allowbreak{}enough-{}\allowbreak{}affine-{}\allowbreak{}weakly-{}\allowbreak{}contractible-{}\allowbreak{}objects})
\item[Proposition 2.23] Exact products, in Ⅰ The six operations (\bpslug{proposition-{}\allowbreak{}exact-{}\allowbreak{}products})
\item[Proposition 2.24] The Milnor sequence, in Ⅰ The six operations (\bpslug{proposition-{}\allowbreak{}milnor-{}\allowbreak{}sequence})
\item[Proposition 2.26] Left completeness, in Ⅰ The six operations (\bpslug{proposition-{}\allowbreak{}left-{}\allowbreak{}completeness})
\item[Corollary 2.29] Continuity of cohomology, in Ⅰ The six operations (\bpslug{corollary-{}\allowbreak{}continuity-{}\allowbreak{}cohomology})
\item[Corollary 2.44] Finite Tor amplitude under comparison, in Ⅰ The six operations (\bpslug{corollary-{}\allowbreak{}finite-{}\allowbreak{}tor-{}\allowbreak{}amplitude-{}\allowbreak{}under-{}\allowbreak{}comparison})
\item[Definition 2.45] Cohomological dimension, in Ⅰ The six operations (\bpslug{definition-{}\allowbreak{}cohomological-{}\allowbreak{}dimension})
\item[Lemma 2.46] Cohomological dimension of the base field, in Ⅰ The six operations (\bpslug{lemma-{}\allowbreak{}cohomological-{}\allowbreak{}dimension-{}\allowbreak{}base-{}\allowbreak{}field})
\item[Proposition 2.49] Unbounded comparison for finite coefficients, in Ⅰ The six operations (\bpslug{proposition-{}\allowbreak{}unbounded-{}\allowbreak{}comparison-{}\allowbreak{}finite-{}\allowbreak{}coefficients})
\item[Example 2.57] The first cohomology of a finite field, in Ⅰ The six operations (\bpslug{example-{}\allowbreak{}first-{}\allowbreak{}cohomology-{}\allowbreak{}finite-{}\allowbreak{}field})
\item[Definition 3.2] Lisse sheaves, in Ⅰ The six operations (\bpslug{definition-{}\allowbreak{}lisse-{}\allowbreak{}sheaves})
\item[Proposition 3.10] Local structure of integral constructible complexes, in Ⅰ The six operations (\bpslug{proposition-{}\allowbreak{}local-{}\allowbreak{}structure-{}\allowbreak{}integral-{}\allowbreak{}constructible-{}\allowbreak{}complexes})
\item[Proposition 3.11] The integral constructibility criterion, in Ⅰ The six operations (\bpslug{proposition-{}\allowbreak{}integral-{}\allowbreak{}constructibility-{}\allowbreak{}criterion})
\item[Section 3.4] Finite Tor amplitude, in Ⅰ The six operations (\bpslug{sec-{}\allowbreak{}Finite-{}\allowbreak{}Tor-{}\allowbreak{}amplitude})
\item[Proposition 3.12] Constructible complexes have finite Tor amplitude for $\ell$-adic coefficients, in Ⅰ The six operations (\bpslug{proposition-{}\allowbreak{}ctf-{}\allowbreak{}equals-{}\allowbreak{}bc})
\item[Lemma 3.13] Flat models and generators, in Ⅰ The six operations (\bpslug{lemma-{}\allowbreak{}flat-{}\allowbreak{}models-{}\allowbreak{}generators})
\item[Lemma 3.17] Descent of lisse $\Qbar$-sheaves, in Ⅰ The six operations (\bpslug{lemma-{}\allowbreak{}descent-{}\allowbreak{}lisse-{}\allowbreak{}qbar-{}\allowbreak{}sheaves-{}\allowbreak{}finite})
\item[Lemma 3.18] Abelian category of constructible sheaves, in Ⅰ The six operations (\bpslug{lemma-{}\allowbreak{}constructible-{}\allowbreak{}sheaves})
\item[Section 3.6] Continuity of Ext, in Ⅰ The six operations (\bpslug{sec-{}\allowbreak{}Continuity-{}\allowbreak{}of-{}\allowbreak{}Ext})
\item[Lemma 3.21] Sections with closed support, in Ⅰ The six operations (\bpslug{lemma-{}\allowbreak{}sections-{}\allowbreak{}support-{}\allowbreak{}closed-{}\allowbreak{}subscheme})
\item[Lemma 3.22] Continuity of Ext, in Ⅰ The six operations (\bpslug{lemma-{}\allowbreak{}continuity-{}\allowbreak{}ext-{}\allowbreak{}out-{}\allowbreak{}constructible-{}\allowbreak{}complex})
\item[Lemma 3.24] Lattices after a finite stratification, in Ⅰ The six operations (\bpslug{lemma-{}\allowbreak{}lattices-{}\allowbreak{}after-{}\allowbreak{}finite-{}\allowbreak{}stratification})
\item[Proposition 4.35] Proper direct image for finite coefficients, in Ⅰ The six operations (\bpslug{proposition-{}\allowbreak{}proper-{}\allowbreak{}direct-{}\allowbreak{}image-{}\allowbreak{}finite-{}\allowbreak{}coefficients})
\item[Proposition 4.40] Proper direct image of constructible complexes, in Ⅰ The six operations (\bpslug{proposition-{}\allowbreak{}proper-{}\allowbreak{}direct-{}\allowbreak{}image-{}\allowbreak{}constructible-{}\allowbreak{}complexes})
\item[Chapter 6] Constructibility of $f^!$, in Ⅰ The six operations (\bpslug{chap-{}\allowbreak{}Constructibility-{}\allowbreak{}of-{}\allowbreak{}f})
\item[Section 6.1] Conventions, in Ⅰ The six operations (\bpslug{sec-{}\allowbreak{}Conventions-{}\allowbreak{}2})
\item[Proposition 6.9] The exceptional inverse image of a smooth morphism, in Ⅰ The six operations (\bpslug{proposition-{}\allowbreak{}exceptional-{}\allowbreak{}inverse-{}\allowbreak{}image-{}\allowbreak{}smooth-{}\allowbreak{}morphism})
\item[Theorem 6.20] $f^!$ with finite coefficients, in Ⅰ The six operations (\bpslug{theorem-{}\allowbreak{}f-{}\allowbreak{}shriek-{}\allowbreak{}finite-{}\allowbreak{}coefficients})
\item[Theorem 6.27] $f^!$ with $\ell$-adic coefficients, in Ⅰ The six operations (\bpslug{theorem-{}\allowbreak{}f-{}\allowbreak{}shriek-{}\allowbreak{}ell-{}\allowbreak{}adic-{}\allowbreak{}coefficients})
\item[Chapter 9] Smooth purity, in Ⅱ Duality and the trace formula (\bpslug{chap-{}\allowbreak{}Smooth-{}\allowbreak{}purity})
\item[Theorem 9.3] Smooth purity for all admissible coefficients, in Ⅱ Duality and the trace formula (\bpslug{theorem-{}\allowbreak{}smooth-{}\allowbreak{}purity-{}\allowbreak{}all-{}\allowbreak{}admissible-{}\allowbreak{}coefficients})
\item[Chapter 10] Verdier and Poincaré duality, in Ⅱ Duality and the trace formula (\bpslug{chap-{}\allowbreak{}Verdier-{}\allowbreak{}and-{}\allowbreak{}Poincare-{}\allowbreak{}duality})
\item[Chapter 12] The trace formula, in Ⅱ Duality and the trace formula (\bpslug{chap-{}\allowbreak{}The-{}\allowbreak{}trace-{}\allowbreak{}formula})
\item[Corollary 12.26] Trace functions and direct image with compact support, in Ⅱ Duality and the trace formula (\bpslug{corollary-{}\allowbreak{}trace-{}\allowbreak{}functions-{}\allowbreak{}direct-{}\allowbreak{}image-{}\allowbreak{}compact})
\item[Definition A.1] Compactification and refinement, in Appendices (\bpslug{definition-{}\allowbreak{}compactification-{}\allowbreak{}refinement})
\item[Theorem A.2] Nagata's compactification theorem, in Appendices (\bpslug{theorem-{}\allowbreak{}nagata-{}\allowbreak{}s-{}\allowbreak{}compactification-{}\allowbreak{}theorem})
\item[Lemma A.5] Passing to the schematic closure, in Appendices (\bpslug{lemma-{}\allowbreak{}passing-{}\allowbreak{}schematic-{}\allowbreak{}closure})
\item[Proposition A.7] Dense compactifications form a cofiltered category, in Appendices (\bpslug{proposition-{}\allowbreak{}dense-{}\allowbreak{}compactifications-{}\allowbreak{}form-{}\allowbreak{}cofiltered-{}\allowbreak{}category})
\end{description}
\clearpage\phantomsection%
\section*{References}
\addcontentsline{toc}{section}{References}
\bpbibitem{BBD82}{bp107}{A. Beilinson, J. Bernstein, and P. Deligne. \emph{Faisceaux pervers}, volume 100 of \emph{Astérisque}. 1982.}

\bpbibitem{BS15}{bp108}{B. Bhatt and P. Scholze. The pro-étale topology for schemes. \emph{Astérisque}, 369:99–201, 2015. \href{https://people.mpim-bonn.mpg.de/scholze/proetale.pdf}{Author's manuscript}.}

\bpbibitem{Con07}{bp109}{B. Conrad. Deligne's notes on Nagata compactifications. \emph{Journal of the Ramanujan Mathematical Society}, 22:205–257, 2007. \href{https://math.stanford.edu/\~conrad/papers/nagatafinal.pdf}{Author's manuscript}.}

\bpbibitem{Del74}{bp110}{P. Deligne. La conjecture de Weil. I. \emph{Publ. Math. IHÉS}, 43:273–307, 1974.}

\bpbibitem{Del80}{bp111}{P. Deligne. La conjecture de Weil. II. \emph{Publ. Math. IHÉS}, 52:137–252, 1980.}

\bpbibitem{Den22}{bp112}{C. Deninger. There is no “Weil-”cohomology theory with real coefficients for arithmetic curves. \emph{Ann. Sc. Norm. Super. Pisa Cl. Sci.}, 2022. To appear. arXiv:\href{https://arxiv.org/abs/2204.02714}{2204.02714}.}

\bpbibitem{Deu41}{bp113}{M. Deuring. Die Typen der Multiplikatorenringe elliptischer Funktionenkörper. \emph{Abh. Math. Sem. Univ. Hamburg}, 14:197–272, 1941.}

\bpbibitem{Eke90}{bp114}{T. Ekedahl. On the adic formalism. In \emph{The Grothendieck Festschrift, Vol. II}, volume 87 of \emph{Progress in Mathematics}, pages 197–218. Birkhäuser, 1990.}

\bpbibitem{Has36}{bp115}{H. Hasse. Zur Theorie der abstrakten elliptischen Funktionenkörper. I, II, III. \emph{J. Reine Angew. Math.}, 175:55–62, 69–88, 193–208, 1936.}

\bpbibitem{Hum72}{bp116}{J. E. Humphreys. \emph{Introduction to Lie Algebras and Representation Theory}, volume 9 of \emph{Graduate Texts in Mathematics}. Springer, 1972.}

\bpbibitem{Ill94}{bp117}{L. Illusie. Autour du théorème de monodromie locale. \emph{Astérisque}, 223:9–57, 1994. §4. \href{https://www.numdam.org/item/AST_1994__223__9_0/}{Original article}.}

\bpbibitem{Kat88}{bp118}{N. M. Katz. Travaux de Laumon. \emph{Astérisque}, 161–162:105–132, 1988. Séminaire Bourbaki, exp. 691.}

\bpbibitem{KW01}{bp119}{R. Kiehl and R. Weissauer. \emph{Weil conjectures, perverse sheaves and \(\ell\)-adic Fourier transform}, volume 42 of \emph{Ergebnisse der Mathematik}. Springer, 2001.}

\bpbibitem{Lau87}{bp120}{G. Laumon. Transformation de Fourier, constantes d'équations fonctionnelles et conjecture de Weil. \emph{Publ. Math. IHÉS}, 65:131–210, 1987.}

\bpbibitem{Mil86a}{bp121}{J. S. Milne. \href{https://www.jmilne.org/math/xnotes/AVs.pdf}{Abelian Varieties}. In \emph{Arithmetic Geometry}. 1986. Corrected online version.}

\bpbibitem{Mil86b}{bp122}{J. S. Milne. \href{https://www.jmilne.org/math/xnotes/JVs.pdf}{Jacobian Varieties}. In \emph{Arithmetic Geometry}. 1986. Corrected online version.}

\bpbibitem{Ran39}{bp123}{R. A. Rankin. Contributions to the theory of Ramanujan's function \(\tau(n)\) and similar arithmetical functions. II. \emph{Proc. Cambridge Philos. Soc.}, 35:357–372, 1939.}

\bpbibitem{SGA4}{bp124}{M. Artin, A. Grothendieck, and J.-L. Verdier, editors. \emph{Théorie des topos et cohomologie étale des schémas}, volumes 269, 270, 305 of \emph{Lecture Notes in Mathematics}. Springer, 1972–1973. SGA 4.}

\bpbibitem{SGA4½}{bp125}{P. Deligne. \emph{Cohomologie étale}, volume 569 of \emph{Lecture Notes in Mathematics}. Springer, 1977. SGA 4½; \emph{Théorèmes de finitude}, §3, and \emph{Rapport sur la formule des traces}, §3, and \emph{Sommes trigonométriques}, §1.}

\bpbibitem{SGA5}{bp126}{A. Grothendieck, editor. \emph{Cohomologie \(\ell\)-adique et fonctions \(L\)}, volume 589 of \emph{Lecture Notes in Mathematics}. Springer, 1977. SGA 5.}

\bpbibitem{SGA7II}{bp127}{P. Deligne and N. Katz, editors. \emph{Groupes de monodromie en géométrie algébrique}, volume 340 of \emph{Lecture Notes in Mathematics}. Springer, 1973. SGA 7 II; Exposé XIII, (2.1.7).}

\bpbibitem{Stacks}{bp128}{The Stacks Project Authors. \href{https://stacks.math.columbia.edu}{The Stacks Project}. Individual results are cited by their stable tags in the text.}

\bpbibitem{Var05}{bp129}{Y. Varshavsky. Lefschetz–Verdier trace formula and a generalization of a theorem of Fujiwara. 2005. Version 2. arXiv:\href{https://arxiv.org/abs/math/0505564}{math/0505564}. \href{https://arxiv.org/abs/math/0505564v2}{Author's preprint}.}

\bpbibitem{Wei48a}{bp130}{A. Weil. \emph{Sur les courbes algébriques et les variétés qui s'en déduisent}. Hermann, Paris, 1948.}

\bpbibitem{Wei48b}{bp131}{A. Weil. \emph{Variétés abéliennes et courbes algébriques}. Hermann, Paris, 1948.}

\bpbibitem{Wei49}{bp132}{A. Weil. Numbers of solutions of equations in finite fields. \emph{Bull. Amer. Math. Soc.}, 55:497–508, 1949.}

\end{document}
